\chapter{Podstawowe rodziny rozkładów prawdopodobieństwa}
\label{ch:distributions}

\providecommand{\atlasref}[2]{#2}

\section{Rozkład Bernoulliego i dwumianowy}
Podstawowym modelem losowym jest doświadczenie z dwiema
rozłącznymi i wyczerpującymi możliwościami: sukces albo porażka. Nie ma potrzeby zakładać,
że są równie prawdopodobne. Parametr $p$ określa szansę sukcesu,
a funkcja charakterystyczna zdarzenia sukcesu ma rozkład Bernoulliego. Licząc sukcesy
w $n$ niezależnych próbach o tym samym parametrze, otrzymujemy
rozkład dwumianowy. Współczynnik $\binom nk$ odpowiada
liczbie różnych rozmieszczeń sukcesów, nie nowej losowości.

\begin{definicja}
Niech $p\in[0,1]$. Mówimy, że zmienna losowa $X$ ma
\emph{rozkład Bernoulliego} z parametrem $p$, jeżeli
$\Prob(X=1)=p$ i $\Prob(X=0)=1-p$.
Piszemy $X\sim\operatorname{Bern}(p)$.
\end{definicja}

Zmienna o rozkładzie Bernoulliego jest prawie na pewno
równa funkcji charakterystycznej $\chi_A$, gdzie $A=\{X=1\}$:
zdarzenia $\{X=0\}$ i $\{X=1\}$ są rozłączne, a suma ich
prawdopodobieństw wynosi jeden. Z własności momentów
funkcji charakterystycznej zdarzenia (rozdział~\ref{ch:expectation})
wynika $\E X=p$ i $\Var(X)=p(1-p)$.
Dla każdego $t\in\R$ ze wzoru \eqref{eq:lotus} dostajemy
$M_X(t)=e^{t\cdot0}(1-p)+e^tp=1-p+pe^t$.
Przy $p=0$ zmienna jest prawie na pewno równa zeru,
a przy $p=1$ --- jedności.

\begin{definicja}
Niech $n\in\N\cup\{0\}$ i $p\in[0,1]$.
Zmienna losowa $X$ ma \emph{rozkład dwumianowy} $\Bin(n,p)$,
gdy dla $k=0,\ldots,n$ zachodzi
\begin{equation}\label{eq:binomial}
\Prob(X=k)=\binom nk p^k(1-p)^{n-k},
\end{equation}
oraz $\Prob(X\in\{0,\ldots,n\})=1$.
Przyjmujemy konwencję $0^0=1$ we wzorze na masy.
Wszystkie masy są nieujemne, a ich suma jest równa
$(p+(1-p))^n=1$ z dwumianu Newtona (również dla $n=0$).
Dla $n=0$ rozkład jest skupiony w zerze, niezależnie od $p$.
Dla $p=0$ jest skupiony w zerze, a dla $p=1$ w punkcie $n$.
Dla $n\geq1$ istnienie niezależnych prób o tym parametrze
wynika z konstrukcji miary produktowej w rozdziale~\ref{ch:probability}.
Dla $n=0$ mamy $\E X=\Var(X)=0$ oraz $M_X(t)=1$
dla każdego $t\in\R$, zgodnie z wzorami z następnego
twierdzenia.
\end{definicja}

\begin{twierdzenie}[Suma prób Bernoulliego]
Niech $n\geq1$ będzie liczbą całkowitą, $p\in[0,1]$,
a $X_1,\ldots,X_n$ niezależnymi zmiennymi losowymi
o rozkładzie $\operatorname{Bern}(p)$ na wspólnej
przestrzeni probabilistycznej. Wtedy
$S_n:=X_1+\cdots+X_n\sim\Bin(n,p)$ oraz
\[
\E S_n=np,\qquad \Var(S_n)=np(1-p),\qquad
M_{S_n}(t)=(1-p+pe^t)^n,\qquad t\in\R.
\]
\end{twierdzenie}
\begin{proof}
Dla ustalonego $k\in\{0,\ldots,n\}$ określmy,
które spośród $n$ prób zakończyły się sukcesem.
Jeśli dokładnie zbiór indeksów $I\subseteq\{1,\ldots,n\}$,
$|I|=k$, odpowiada sukcesom, to zdarzenie to ma postać
$\bigcap_{i\in I}\{X_i=1\}\cap
\bigcap_{j\notin I}\{X_j=0\}$.
Z niezależności jego prawdopodobieństwo jest równe
$p^k(1-p)^{n-k}$.

Przekrój po pustym zbiorze indeksów oznacza całą
przestrzeń probabilistyczną, a pusty iloczyn wynosi jeden.
Różne zbiory $I$ opisują rozłączne zdarzenia,
a ich suma różni się od zdarzenia $\{S_n=k\}$ co najwyżej
o podzbiór mierzalnego zbioru zerowego
$\bigcup_{i=1}^n\{X_i\notin\{0,1\}\}$.
W konsekwencji mają to samo prawdopodobieństwo. Ponieważ
liczba $k$-elementowych podzbiorów indeksów wynosi
$\binom nk$, z addytywności dostajemy
$\Prob(S_n=k)=\binom nk p^k(1-p)^{n-k}$.
Po zsumowaniu tego wyrażenia dla wszystkich $k$
otrzymujemy przez dwumian Newtona
$(p+(1-p))^n=1$, co stanowi dodatkową kontrolę normalizacji.

Dla każdej pojedynczej próby mamy
$\E X_i=p$, $\Var(X_i)=p(1-p)$.
Wszystkie zmienne $X_i$ mają skończone drugie momenty,
ponieważ $|X_i|\leq1$ prawie na pewno. Suma również
ma skończony drugi moment, gdyż $|S_n|\leq n$ prawie na pewno.
Liniowość wartości oczekiwanej i addytywność
wariancji niezależnych składników (rozdział~\ref{ch:expectation}) dają
$\E S_n=np$ i $\Var(S_n)=np(1-p)$.
Wreszcie $M_{X_i}(t)=1-p+pe^t$ jest skończona dla każdego
$t\in\R$. Stosując wzór na funkcję tworzącą momenty sumy
niezależnych zmiennych z rozdziału~\ref{ch:expectation}
kolejno do sum częściowych (każda suma poprzednich zmiennych
jest niezależna od następnej na mocy produktowości rozkładu
łącznego), otrzymujemy
$M_{S_n}(t)=\prod_{i=1}^nM_{X_i}(t)
=(1-p+pe^t)^n$.
\end{proof}

\begin{przyklad}[Liczba wadliwych elementów]
Jeżeli każdy z 20 elementów jest wadliwy niezależnie z
prawdopodobieństwem $0{,}02$, to liczba wadliwych elementów
ma rozkład $\Bin(20,0{,}02)$. Prawdopodobieństwo
braku wad wynosi $0{,}98^{20}$, a dokładnie jednej wady
$20\cdot0{,}02\cdot0{,}98^{19}$.
Jeżeli elementy pobierano z małej partii \emph{bez zwracania},
rozkład dwumianowy nie jest na ogół dokładny.
\end{przyklad}

\section{Rozkład geometryczny i ujemny dwumianowy}
Rozkład dwumianowy ustala liczbę prób i pyta o liczbę sukcesów.
Można odwrócić punkt widzenia: ustalić liczbę sukcesów i
pytać, ile prób potrzeba, by je osiągnąć. Pierwszy sukces
prowadzi do rozkładu geometrycznego, a $r$-ty do ujemnego
dwumianowego. Ich postać odzwierciedla zdarzenie
„najpierw odpowiednia liczba porażek, a następnie sukces”.

\begin{definicja}
Niech $p\in(0,1]$. Rozkładem geometrycznym $\Geom(p)$
na $\{1,2,\ldots\}$ nazywamy rozkład o funkcji masy
$\Prob(T=k)=(1-p)^{k-1}p$, $k\geq1$ całkowity,
przy konwencji $0^0=1$.
Masy sumują się do jedności: przy $q=1-p\in[0,1)$
mamy $\sum_{k=1}^\infty pq^{k-1}=p/(1-q)=1$.
Dla $p=1$ rozkład jest skupiony w punkcie $1$.

W modelu niezależnych prób Bernoulliego (którego istnienie
wykazano w rozdziale~\ref{ch:probability}) niech $T$ będzie
numerem pierwszej udanej próby, z wartością $\infty$,
gdy sukces nigdy nie nastąpi. Możemy przyjąć przestrzeń
ciągów o współrzędnych w $\{0,1\}$ z miarą produktową,
tak że każda próba ma dokładnie dwie możliwe wartości
w każdym punkcie przestrzeni. Zdarzenie $\{T=k\}$ jest
przekrojem zdarzeń $k-1$ porażek i jednego sukcesu,
więc jest mierzalne i ma prawdopodobieństwo $q^{k-1}p$.
Zdarzenie $\{T=\infty\}$ jest przekrojem malejącego
ciągu zdarzeń $C_k$, że pierwszych $k$ prób zakończyło się
porażkami. Z niezależności $\Prob(C_k)=q^k$, a z ciągłości
miary z góry (rozdział~\ref{ch:probability}) wynika
$\Prob(T=\infty)=\lim_{k\to\infty}q^k=0$. Zastąpienie wartości $\infty$
przez $1$ na tym mierzalnym zbiorze zerowym daje rzeczywistą
zmienną losową o rozkładzie $\Geom(p)$: przeciwobrazy
pojedynczych wartości są mierzalne, a przeciwobraz
dowolnego zbioru borelowskiego jest przeliczalną sumą
takich przeciwobrazów.
\end{definicja}

\begin{stwierdzenie}
Dla $T\sim\Geom(p)$ zachodzi $\Prob(T>k)=(1-p)^k$ dla
$k\in\N\cup\{0\}$, a także
$\E T=1/p$ i $\Var(T)=(1-p)/p^2$.
Jeśli $p\in(0,1)$, to dla całkowitych $m,n\geq0$
$\Prob(T>m+n\mid T>m)=\Prob(T>n)$.
\end{stwierdzenie}
\begin{proof}
Z definicji rozkładu dla całkowitego $k\geq0$ mamy
$\Prob(T>k)=\sum_{j=k+1}^\infty pq^{j-1}=q^k$,
gdzie $q=1-p$ (dla $q=0$ i $k=0$ stosujemy $0^0=1$).
Wzór obowiązuje dla każdej zmiennej o tym rozkładzie,
nie tylko dla zmiennej skonstruowanej z prób.

Dla $q<1$ szereg geometryczny
$\sum_{k=0}^\infty q^k=1/(1-q)$ jest zbieżny
jednostajnie na każdym przedziale $[0,q_0]$ z $q_0<1$.
Również szeregi pochodnych, począwszy od
$\sum_{k=1}^\infty kq^{k-1}$, są jednostajnie zbieżne
na takich przedziałach dla $0<q_0<1$.
Dla pierwszej pochodnej stosujemy majorantę $kq_0^{k-1}$,
a dla drugiej $k(k-1)q_0^{k-2}$, $k\geq2$;
oba szeregi majorant są zbieżne z kryterium d'Alemberta.
Twierdzenie o różniczkowaniu szeregów potęgowych
pozwala różniczkować wyraz po wyrazie:
\[
\sum_{k=1}^\infty kq^{k-1}=\frac1{(1-q)^2},
\qquad
\sum_{k=2}^\infty k(k-1)q^{k-2}=\frac2{(1-q)^3}.
\]
Z tożsamości $k^2=k(k-1)+k$ wynika więc
\[
\sum_{k=1}^\infty k^2q^{k-1}
=\frac{2q}{(1-q)^3}+\frac1{(1-q)^2}
=\frac{1+q}{(1-q)^3}.
\]
Wzory obowiązują również dla $q=0$.
Ze wzoru \eqref{eq:lotus} dla nieujemnych funkcji
$x\mapsto |x|$ i $x\mapsto x^2$, korzystając z faktu,
że $T\geq1$ prawie na pewno, otrzymujemy skończone momenty:
$\E T=p/(1-q)^2=1/p$,
$\E T^2=p(1+q)/(1-q)^3=(1+q)/p^2$
i $\Var(T)=\E T^2-(\E T)^2=q/p^2$.

Dla $p\in(0,1)$ mamy $q>0$ i
$\Prob(T>m)>0$, więc warunkowanie jest zdefiniowane.
Ponieważ $\{T>m+n\}\subseteq\{T>m\}$,
\[
\Prob(T>m+n\mid T>m)
=\frac{q^{m+n}}{q^m}=q^n=\Prob(T>n).
\]
Przy $p=1$ zdarzenie $\{T>m\}$ dla $m\geq1$
ma masę zero i powyższe warunkowanie nie ma sensu,
co wyjaśnia ograniczenie w tezie. Dla $p=1$ i $m=0$
warunkowanie jest określone, ponieważ $\Prob(T>0)=1$,
a podana równość zachodzi bezpośrednio.
\end{proof}

\begin{uwaga}
Spotyka się także rozkład geometryczny na $\{0,1,\ldots\}$,
gdzie zmienna oznacza liczbę niepowodzeń \emph{przed} pierwszym sukcesem.
Jest to rozkład $T-1$, o masach $pq^k$, $k\geq0$,
gdzie $q=1-p$. Z algebry momentów wynika, że jego średnia
wynosi $1/p-1=(1-p)/p$, a wariancja pozostaje równa $q/p^2$.
Dlatego zawsze należy podać konwencję numerowania.
\end{uwaga}

\begin{definicja}
Dla całkowitego $r\geq1$ i $p\in(0,1]$ rozkładem ujemnym
dwumianowym (w konwencji liczby prób do $r$-tego sukcesu)
nazywamy rozkład na $\{r,r+1,\ldots\}$ o masach
$\Prob(T_r=n)=\binom{n-1}{r-1}p^r(1-p)^{n-r}$.
Przy $p=1$ przyjmujemy $0^0=1$, więc rozkład jest
skupiony w $r$. Normalizację dla wszystkich parametrów
uzasadnia konstrukcja w następnym dowodzie.
\end{definicja}

\begin{stwierdzenie}
Jeżeli $T_r$ jest numerem próby, w której następuje $r$-ty sukces
w ciągu niezależnych prób Bernoulliego na wspólnej przestrzeni
probabilistycznej z parametrem
$p\in(0,1]$, przy całkowitym $r\geq1$, to jest skończony
prawie na pewno i ma powyższy rozkład (po dowolnym
skończonym mierzalnym przedłużeniu na zbiorze zerowym),
a ponadto
$\E T_r=r/p$ i $\Var(T_r)=r(1-p)/p^2$.
\end{stwierdzenie}
\begin{proof}
Najpierw definiujemy $T_r=\infty$, jeśli $r$-ty sukces
nigdy nie nastąpi. Możemy traktować próby jako zmienne
przyjmujące wyłącznie wartości $0,1$: jeśli pierwotne
zmienne mają te wartości jedynie prawie na pewno,
modyfikujemy je na mierzalnej sumie przeliczalnie wielu
zbiorów zerowych, przyjmując tam wartości $0$.
Rozkłady wszystkich skończonych wektorów prób,
a więc również ich niezależność, pozostają niezmienione. Każde zdarzenie $\{T_r=n\}$ jest
mierzalne, ponieważ zależy od skończonej liczby prób.
Zdarzenia sukcesu w kolejnych próbach są wzajemnie
niezależne, a suma ich prawdopodobieństw jest nieskończona,
gdyż $p>0$. Z drugiego lematu Borela--Cantellego
(rozdział~\ref{ch:probability}) prawie na pewno występuje
nieskończenie wiele sukcesów. W szczególności $T_r$ jest
skończony prawie na pewno i można przyjąć $T_r=r$ na
mierzalnym zdarzeniu $\{T_r=\infty\}$.

Aby $r$-ty sukces nastąpił dokładnie w próbie numer $n$,
próba $n$ musi zakończyć się sukcesem, a wśród pierwszych
$n-1$ prób musi znaleźć się dokładnie $r-1$ sukcesów.
Dla każdej konkretnej konfiguracji tych prób, przez
niezależność, prawdopodobieństwo wynosi
$p^r(1-p)^{n-r}$. Konfiguracji jest
$\binom{n-1}{r-1}$ (rozdział~\ref{ch:combinatorics});
są rozłączne, co daje wzór na masę. Ewentualna modyfikacja
na zbiorze zerowym nie zmienia tych prawdopodobieństw.
Ponieważ $T_r$ jest skończony prawie na pewno i przyjmuje
wartości co najmniej $r$, suma obliczonych mas wynosi jeden.

Wykażmy niezależność czasów oczekiwania.
Dla $1\leq j\leq r$ niech $T_j$ oznacza numer próby
$j$-tego sukcesu, określony analogicznie do $T_r$.
Poza zdarzeniem, na którym nie występuje $r$ sukcesów,
wszystkie te liczby są skończone.
Niech $W_1$ będzie liczbą prób do pierwszego sukcesu,
a dla $j\geq2$ niech $W_j$ oznacza liczbę kolejnych prób
pomiędzy sukcesami numer $j-1$ i $j$, licząc próbę
zakończoną sukcesem. Na mierzalnym zbiorze zerowym,
na którym nie występuje $r$ sukcesów, definiujemy wszystkie
$W_j=1$. Zdarzenia określające wartości $W_j$ są mierzalne:
dla $j\geq2$, poza tym zbiorem, zdarzenie $\{W_j=k\}$
jest sumą po $m\geq j-1$ zdarzeń
$\{T_{j-1}=m,T_j=m+k\}$; dla $j=1$ korzystamy
z mierzalności $\{T_1=k\}$. Uwzględnienie mierzalnego
zbioru zerowego, na którym $W_j=1$, nie narusza mierzalności.
Dla dowolnych całkowitych $k_1,\ldots,k_r\geq1$
zdarzenie $\{W_1=k_1,\ldots,W_r=k_r\}$ oznacza, z dokładnością do zbioru zerowego,
jednoznacznie określony ciąg
$k_1+\cdots+k_r$ wyników: w każdym bloku $k_j-1$
niepowodzeń, po których następuje sukces.
Przez niezależność pojedynczych prób jego
prawdopodobieństwo jest równe
$\prod_{j=1}^r p(1-p)^{k_j-1}$.
Jest to iloczyn mas rozkładów $\Geom(p)$.
Sumując tę równość po dowolnych zbiorach wartości
$k_j\in A_j\subseteq\{1,2,\ldots\}$, otrzymujemy
\[
\Prob(W_j\in A_j,\ j=1,\ldots,r)
=\prod_{j=1}^r\sum_{k\in A_j}p(1-p)^{k-1}.
\]
Rozdzielenie sum jest uprawnione przez nieujemność wyrazów.
Ustalając jeden zbiór $A_j$ i biorąc wszystkie pozostałe
równe $\{1,2,\ldots\}$, otrzymujemy rozkłady brzegowe
$\Geom(p)$, a następnie z powyższej faktoryzacji
niezależność. Wszystkie $W_j$ mają skończone drugie
momenty, zatem $T_r=W_1+\cdots+W_r$ prawie na pewno
również ma skończony drugi moment, na mocy nierówności
$(\sum_{j=1}^r W_j)^2\leq r\sum_{j=1}^r W_j^2$.
Z liniowości wartości oczekiwanej i addytywności
wariancji niezależnych składników otrzymujemy
$\E T_r=r/p$ i $\Var(T_r)=r(1-p)/p^2$.
\end{proof}

\section{Rozkład hipergeometryczny}
Liczba sukcesów w próbie bez zwracania różni się od liczby
sukcesów w niezależnych próbach Bernoulliego. Każdy kolejny
wynik zmienia skład dostępnej populacji. Zliczając możliwe
podzbiory zamiast ciągów niezależnych doświadczeń, otrzymujemy
rozkład hipergeometryczny. Niedodatnia kowariancja wyników kolejnych
losowań tłumaczy, dlaczego wariancja nie przekracza wariancji
w modelu dwumianowym z taką samą liczbą prób i średnią.

\begin{definicja}
Niech $N\geq1$ będzie liczbą całkowitą,
$K\in\{0,\ldots,N\}$ i $n\in\{0,\ldots,N\}$.
Zmienna losowa $X$ ma rozkład hipergeometryczny
$\operatorname{HGeom}(N,K,n)$, jeżeli
\begin{equation}\label{eq:hypergeom}
\Prob(X=k)=\frac{\binom Kk\binom{N-K}{n-k}}{\binom Nn}
\quad(k\in\Z),
\end{equation}
przy konwencji, że symbol Newtona jest zerem poza właściwym zakresem,
oraz $\Prob(X\in\mathbb Z)=1$.
Masy są dodatnie dokładnie dla całkowitych
$\max(0,n-(N-K))\leq k\leq\min(n,K)$.
Aby uzasadnić wzór i normalizację, rozważmy wszystkie
$\binom Nn$ podzbiory $n$-elementowe populacji $N$ różnych
elementów, z których $K$ jest wyróżnionych, z jednakowymi
prawdopodobieństwami. Podzbiorów zawierających dokładnie
$k$ wyróżnionych elementów jest $\binom Kk\binom{N-K}{n-k}$,
ponieważ oba wybory wykonujemy z rozłącznych części populacji
(rozdział~\ref{ch:combinatorics}). Podział wszystkich
podzbiorów według liczby wyróżnionych elementów daje
\[
\sum_{k\in\mathbb Z}\binom Kk\binom{N-K}{n-k}=\binom Nn.
\]
Otrzymujemy zatem rozkład prawdopodobieństwa i jego
interpretację jako liczby wyróżnionych obiektów w próbie
losowanej jednostajnie bez zwracania. Ponadto $0\leq X\leq n$
prawie na pewno, więc wszystkie jego dodatnie momenty
bezwzględne są skończone.
\end{definicja}

\begin{twierdzenie}
Jeżeli $X\sim\operatorname{HGeom}(N,K,n)$, to
$\E X=nK/N$. Dla $N>1$ zachodzi
\[
\Var(X)=n\frac KN\left(1-\frac KN\right)\frac{N-n}{N-1}.
\]
Dla $N=1$ wariancja wynosi zero.
\end{twierdzenie}
\begin{proof}
Wystarczy obliczyć momenty dowolnej zmiennej o zadanym
rozkładzie, gdyż zależą one wyłącznie od rozkładu
(ze wzoru \eqref{eq:lotus}).
Realizujmy losowanie jako pierwszych $n$ pozycji jednostajnej
losowej permutacji $N$ różnych elementów. Każdy ustalony
podzbiór $n$-elementowy występuje na tych pozycjach
w $n!(N-n)!$ spośród $N!$ permutacji, a więc z prawdopodobieństwem
$1/\binom Nn$. Dla $j=1,\ldots,n$ niech $A_j$ będzie
zdarzeniem, że na pozycji $j$ znajduje się obiekt wyróżniony,
oraz $I_j=\chi_{A_j}$. Liczba wyróżnionych elementów
$X=\sum_{j=1}^n I_j$ ma zatem żądany rozkład.
Dla $n=0$ suma jest pusta i równa zeru, zatem oba
wzory na momenty zachodzą bezpośrednio. Dalej przyjmijmy
$n\geq1$. Każdy z $N$ elementów występuje na ustalonej pozycji
w $(N-1)!$ permutacjach, toteż
$\Prob(I_j=1)=K/N=:p$ i $\E I_j=p$.
Liniowość wartości oczekiwanej daje $\E X=np$.

Dla $N>1$ i różnych $i,j$ każda uporządkowana para
różnych elementów występuje na pozycjach $i,j$ w $(N-2)!$
permutacjach. Jeżeli $K\geq2$, par wyróżnionych elementów
jest $K(K-1)$, a jeżeli $K\in\{0,1\}$, nie ma żadnej.
We wszystkich przypadkach otrzymujemy więc
$\Prob(I_i=I_j=1)=K(K-1)/(N(N-1))$.
Ten argument nie wymaga warunkowania względem zdarzenia
o prawdopodobieństwie zero.
Zatem
\[
\Cov(I_i,I_j)=\frac{K(K-1)}{N(N-1)}-
\left(\frac KN\right)^2
=-\frac{p(1-p)}{N-1}.
\]
Każda funkcja charakterystyczna $I_j$ ma wariancję $p(1-p)$.
Rozwijając kwadrat $\sum_{j=1}^n(I_j-p)$ i całkując,
otrzymujemy $n$ wariancji i $n(n-1)$ kowariancji
uporządkowanych par różnych indeksów. Otrzymujemy
\[
\Var(X)=np(1-p)-\frac{n(n-1)}{N-1}p(1-p)
=np(1-p)\frac{N-n}{N-1}.
\]
Jeżeli $N=1$, to albo $n=0$ (przypadek już rozpatrzony),
albo $n=1$, gdy $X=K$ prawie na pewno, więc
$\Var(X)=0$ i $\E X=K=nK/N$.
\end{proof}

\begin{uwaga}
Dla $N>1$ i $1\leq n\leq N$ względem wariancji
rozkładu $\Bin(n,K/N)$ pojawia się czynnik korekcyjny
$(N-n)/(N-1)\in[0,1]$.
Wariancja jest ściśle mniejsza, gdy $n\geq2$ i $0<K<N$.
Dla $n=0$ obie wariancje wynoszą zero.
Gdy $n/N$ jest małe, a $N$ duże, czynnik korekcyjny
jest bliski jedności; samo porównanie wariancji nie
dowodzi jednak bliskości całych rozkładów.
\end{uwaga}


\begin{przyklad}[Taka sama średnia, ale różna wariancja]
W populacji znajduje się $N=10$ kul, z których $K=5$
jest czerwonych. Losujemy $n=4$ kule bez zwracania.
Liczba czerwonych ma rozkład hipergeometryczny,
a więc średnią $4\cdot(5/10)=2$ oraz wariancję
$4\cdot(1/2)(1/2)\cdot(10-4)/(10-1)=2/3$.
Gdybyśmy losowali ze zwracaniem, liczba czerwonych
miałaby rozkład $\Bin(4,1/2)$ o tej samej średniej $2$,
ale wariancji $4\cdot(1/2)(1/2)=1$.
Przy losowaniu bez zwracania pojawienie się czerwonej kuli
zmniejsza szansę czerwonej kuli w kolejnej próbie.
Ta ujemna zależność ogranicza rozrzut liczby sukcesów.
Rysunek~\ref{fig:rp1-06-hypergeom-binomial} przedstawia porównanie obu funkcji masy.
\end{przyklad}

\begin{figure}[htbp]
\centering
\begin{tikzpicture}[x=1.65cm,y=7cm,>=Stealth]
\draw[->] (-.45,0)--(4.65,0) node[right] {$k$};
\draw[->] (-.43,0)--(-.43,.55) node[above] {$\Prob(X=k)$};
\foreach \y/\txt in {.1/{0{,}1},.2/{0{,}2},.3/{0{,}3},.4/{0{,}4},.5/{0{,}5}} {
 \draw[gray!20] (-.43,\y)--(4.45,\y);
 \node[left,font=\scriptsize] at (-.43,\y) {$\txt$};
}
\foreach \k/\h/\b in {0/.02380952/.0625,1/.23809524/.25,2/.47619048/.375,3/.23809524/.25,4/.02380952/.0625} {
 \fill[blue!70!black] (\k-.25,0) rectangle (\k-.035,\h);
 \fill[orange!80!black] (\k+.035,0) rectangle (\k+.25,\b);
 \node[below,font=\small] at (\k,0) {$\k$};
}
\fill[blue!70!black] (.1,-.1) rectangle (.28,-.075);
\node[right,font=\small] at (.32,-.0875) {bez zwracania};
\fill[orange!80!black] (2.3,-.1) rectangle (2.48,-.075);
\node[right,font=\small] at (2.52,-.0875) {ze zwracaniem};
\end{tikzpicture}
\caption{Funkcje masy rozkładu hipergeometrycznego
($N=10$, $K=5$, $n=4$) i dwumianowego ($n=4$, $p=1/2$).
Oba rozkłady mają średnią $2$, lecz pierwszy ma
mniejszą wariancję i większą masę w punkcie $2$.}
\label{fig:rp1-06-hypergeom-binomial}
\end{figure}


\section{Rozkład Poissona}
Jeśli mamy bardzo wiele prób, lecz sukces w każdej z nich jest
bardzo rzadki, rozkład dwumianowy bywa niewygodny w obliczeniach.
Gdy liczba prób rośnie, prawdopodobieństwo sukcesu maleje, a
oczekiwana liczba sukcesów pozostaje stała, otrzymujemy
rozkład Poissona. Parametr $\lambda$ opisuje jednocześnie
średnią i wariancję. Samo występowanie zdarzeń „w czasie” nie
wystarcza jednak do przyjęcia rozkładu Poissona: odpowiednia
struktura zależności musi wynikać z modelu.

\begin{definicja}
Niech $\lambda>0$. Zmienna losowa $X$ ma rozkład Poissona
$\Poi(\lambda)$, jeżeli
$\Prob(X=k)=e^{-\lambda}\lambda^k/k!$ dla
$k\in\{0,1,2,\ldots\}$. Suma tych nieujemnych mas
wynosi jeden (co sprawdzimy w następnym dowodzie),
więc rozkład jest skupiony na tym zbiorze.
\end{definicja}

\begin{stwierdzenie}
Dla $X\sim\Poi(\lambda)$ zachodzi
$\E X=\lambda$, $\Var(X)=\lambda$ oraz
$M_X(t)=\exp(\lambda(e^t-1))$ dla $t\in\R$.
\end{stwierdzenie}
\begin{proof}
Najpierw sprawdzamy, że wzór definiuje rozkład:
\[
\sum_{k=0}^\infty e^{-\lambda}\frac{\lambda^k}{k!}
=e^{-\lambda}e^\lambda=1.
\]
Dla każdego $t\in\R$ stosujemy wzór \eqref{eq:lotus}
do nieujemnej funkcji $x\mapsto e^{tx}$.
Funkcja tworząca momenty jest skończona, ponieważ
\[
M_X(t)=\sum_{k=0}^\infty e^{tk}
 e^{-\lambda}\frac{\lambda^k}{k!}
=e^{-\lambda}\sum_{k=0}^\infty\frac{(\lambda e^t)^k}{k!}
=\exp\bigl(\lambda(e^t-1)\bigr).
\]
Ponieważ funkcja jest skończona w otoczeniu zera,
możemy odzyskać momenty przez różniczkowanie, zgodnie
z twierdzeniem z rozdziału~\ref{ch:expectation}:
$M_X'(t)=\lambda e^tM_X(t)$,
a $M_X''(t)=(\lambda e^t+\lambda^2e^{2t})M_X(t)$.
W szczególności $\E X=M_X'(0)=\lambda$,
$\E X^2=M_X''(0)=\lambda+\lambda^2$,
stąd $\Var(X)=\lambda$.
\end{proof}

\begin{twierdzenie}[Prawo rzadkich zdarzeń]
Niech $\lambda>0$ i niech $X_n\sim\Bin(n,\lambda/n)$
dla całkowitych $n\geq\lceil\lambda\rceil$ z $n\geq1$.
Wtedy dla każdego ustalonego $k\in\N\cup\{0\}$
\[
\lim_{n\to\infty}\Prob(X_n=k)
=e^{-\lambda}\frac{\lambda^k}{k!}.
\]
\end{twierdzenie}
\begin{proof}
Dla ustalonego $k\geq0$ i dostatecznie dużego $n$
mamy $n\geq k$ oraz $\lambda/n<1$. Z definicji
rozkładu dwumianowego
\[
\Prob(X_n=k)=\binom nk\left(\frac\lambda n\right)^k
\left(1-\frac\lambda n\right)^{n-k}.
\]
Rozdzielamy to wyrażenie na czynniki o prostych granicach:
\[
\Prob(X_n=k)=\frac{\lambda^k}{k!}
\underbrace{\prod_{j=0}^{k-1}\left(1-\frac jn\right)}_{\longrightarrow1}
\underbrace{\left(1-\frac\lambda n\right)^n}_{\longrightarrow e^{-\lambda}}
\underbrace{\left(1-\frac\lambda n\right)^{-k}}_{\longrightarrow1}.
\]
Pierwszy iloczyn ma stałą liczbę czynników (bo $k$ nie
zależy od $n$). Granica środkowego czynnika jest
klasyczną granicą wykładniczą. Istotnie, z różniczkowalności
logarytmu w $1$ wynika $\log(1-u)/u\to-1$ przy $u\downarrow0$,
zatem
\[
n\log(1-\lambda/n)=\lambda\frac{\log(1-\lambda/n)}{\lambda/n}
\longrightarrow-\lambda.
\]
Ciągłość funkcji wykładniczej daje żądaną granicę.
Dla $k=0$ pierwszy iloczyn jest pusty i wynosi jeden.
Ostatni czynnik zmierza do jedności.
Iloczyn granic daje $e^{-\lambda}\lambda^k/k!$.
Dla dowolnego $x\geq0$ dystrybuanta $X_n$ jest
skończoną sumą mas dla $k=0,\ldots,\lfloor x\rfloor$,
więc zbiega do analogicznej sumy mas rozkładu Poissona.
Dla $x<0$ obie dystrybuanty są zerowe.
Otrzymujemy zbieżność dystrybuant w każdym punkcie,
a tym samym zbieżność według rozkładu w sensie
definicji z rozdziału~\ref{ch:convergence}.
\end{proof}

\begin{przyklad}[Zliczanie defektów]
Jeśli w dużej liczbie $n$ niezależnych prób każda próba
kończy się defektem z tym samym małym prawdopodobieństwem
$\lambda/n$, to zgodnie z prawem rzadkich zdarzeń
rozkład Poissona przybliża rozkład liczby defektów
(dla ustalonych wartości liczby defektów przy $n\to\infty$).
Nie jest jednak automatycznym modelem dowolnych zliczeń:
przy klastrowaniu zdarzeń lub nadmiernej zmienności
wariancja może różnić się istotnie od średniej.
\end{przyklad}

\begin{stwierdzenie}[Dodawanie rozkładów Poissona]
Jeśli $X\sim\Poi(\lambda)$ i $Y\sim\Poi(\eta)$ są niezależne
na wspólnej przestrzeni probabilistycznej,
gdzie $\lambda,\eta>0$, to $X+Y\sim\Poi(\lambda+\eta)$.
\end{stwierdzenie}
\begin{proof}
Obie zmienne przyjmują prawie na pewno wartości
całkowite nieujemne. Zdarzenie $\{X+Y=k\}$ dla całkowitego
$k\geq0$ jest, z dokładnością do zbioru zerowego,
sumą rozłącznych zdarzeń $\{X=j,Y=k-j\}$, $0\leq j\leq k$.
Z addytywności i niezależności wynika
\[
\Prob(X+Y=k)
=\sum_{j=0}^k e^{-\lambda}\frac{\lambda^j}{j!}
 e^{-\eta}\frac{\eta^{k-j}}{(k-j)!}
=e^{-(\lambda+\eta)}\frac{(\lambda+\eta)^k}{k!},
\]
gdzie ostatnią równość daje dwumian Newtona.
\end{proof}

\begin{przyklad}[Oczekiwanie na zebranie wszystkich kuponów]
Niech $N\geq1$ będzie liczbą całkowitą, a $U_1,U_2,\ldots$
niech będą niezależnymi zmiennymi o rozkładzie jednostajnym
na zbiorze $\{1,\ldots,N\}$. Możemy przyjąć, że wszystkie
zmienne mają wartości w tym zbiorze w każdym punkcie:
ewentualne modyfikacje na mierzalnej sumie zbiorów
zerowych nie zmieniają rozkładów łącznych prób.
Niech $T_j$ oznacza pierwszy numer próby, po której
widziano $j$ różnych etykiet,
przy $T_0=0$ i konwencji $T_j=\infty$, gdy takiej próby nie ma.
Dla każdej etykiety $a$ prawdopodobieństwo, że nie pojawiła
się ona w pierwszych $m\geq1$ próbach, wynosi $(1-1/N)^m$
z niezależności prób. Zdarzenie $\{T_N>m\}$ jest sumą
zdarzeń niepojawienia się poszczególnych etykiet.
Z podaddytywności miary wynika
$\Prob(T_N>m)\leq N(1-1/N)^m\to0$.
Ponieważ $\{T_N=\infty\}=\bigcap_{m\geq1}\{T_N>m\}$,
ciągłość miary z góry daje $\Prob(T_N=\infty)=0$.
Wtedy również wszystkie $T_j$, $0\leq j\leq N$,
są skończone prawie na pewno.
Zdarzenia $\{T_j=m\}$ są mierzalne, ponieważ zależą tylko
od pierwszych $m$ prób.

Dla $0\leq j<N$ oznaczmy $W_j=T_{j+1}-T_j$ poza
mierzalnym zbiorem zerowym $\{T_N=\infty\}$; na nim
przyjmijmy $W_j=1$. Dla każdego całkowitego $k\geq1$
zdarzenie $\{W_j=k\}$ poza zbiorem $\{T_N=\infty\}$
jest sumą po całkowitych $m\geq0$ zdarzeń
$\{T_j=m,T_{j+1}=m+k\}$ ograniczonych do dopełnienia
tego zbioru zerowego. Stąd $W_j$ są mierzalnymi
zmiennymi o wartościach w $\{1,2,\ldots\}$.
Aby wyznaczyć ich rozkład, dla $m\geq0$ i
$A\subseteq\{1,\ldots,N\}$, $|A|=j$, rozważmy
$B_{m,A}=\{T_j=m,\ \{U_1,\ldots,U_m\}=A\}$.
Dla $j=0$ chodzi o $m=0$ i $A=\varnothing$.
Zdarzenia $B_{m,A}$ tworzą, z dokładnością do zbioru
zerowego, przeliczalny rozłączny podział przestrzeni.
Dla całkowitego $k\geq0$ zdarzenie
$B_{m,A}\cap\{W_j>k\}$ jest, z dokładnością do zbioru
zerowego, równe
$B_{m,A}\cap\bigcap_{\ell=1}^k\{U_{m+\ell}\in A\}$.
Zdarzenie $B_{m,A}$ zależy tylko od pierwszych $m$ prób,
a pozostałe zdarzenia dotyczą kolejnych $k$ prób.
Z produktowości skończonych rozkładów łącznych wynika,
że prawdopodobieństwo przekroju jest iloczynem
$\Prob(B_{m,A})(j/N)^k$, z konwencją $0^0=1$.
Dla $k=0$ przekrój pustej rodziny zdarzeń oznacza
całą przestrzeń probabilistyczną.
Sumując po $m,A$, uzyskujemy
$\Prob(W_j>k)=(j/N)^k$.
Odejmując kolejne wartości ogona, dostajemy rozkład
geometryczny $\Geom((N-j)/N)$ oraz
$\E W_j=N/(N-j)$.
Niech $T=T_N$ po zastąpieniu wartości nieskończonej
dowolną skończoną wartością na zbiorze zerowym.
Ponieważ $T=\sum_{j=0}^{N-1}W_j$ prawie na pewno,
z liniowości wartości oczekiwanej otrzymujemy
$\E T=\sum_{j=0}^{N-1}N/(N-j)=N\sum_{k=1}^N1/k$.
Nie trzeba przy tym zakładać, że wszystkie czasy oczekiwania
są z góry wzajemnie niezależne: do obliczenia średniej
wystarcza liniowość wartości oczekiwanej.
\end{przyklad}

\section{Rozkład jednostajny ciągły i wykładniczy}
Nie można przypisać jednakowej dodatniej masy wszystkim
punktom niejednopunktowego odcinka: już suma mas dostatecznie
wielu różnych punktów przekroczyłaby jeden.
Dla rozkładu jednostajnego na $(a,b)$ prawdopodobieństwo
przedziału zawartego w $(a,b)$ jest proporcjonalne do jego
długości. Rozkład wykładniczy odpowiada innej intuicji:
jeśli wydarzenie nie zaszło przez czas $s$, dodatkowy czas
oczekiwania ma mieć tę samą regułę co na początku.
Poniżej sprawdzimy, że rozkład wykładniczy ma właśnie
taką własność braku pamięci.

\begin{definicja}
Dla $a,b\in\R$ z $a<b$ rozkładem jednostajnym
$\Unif(a,b)$ nazywamy rozkład o gęstości
$f(x)=\chi_{(a,b)}(x)/(b-a)$ względem miary Lebesgue'a
$\lambda_1$. Wybór wartości gęstości w punktach $a,b$
nie zmienia rozkładu.
\end{definicja}

\begin{stwierdzenie}
Jeżeli $X\sim\Unif(a,b)$, to
\[
F_X(x)=\begin{cases}
0,&x\leq a,\\
\dfrac{x-a}{b-a},&a<x<b,\\
1,&x\geq b,
\end{cases}
\qquad
\E X=\frac{a+b}{2},\qquad
\Var(X)=\frac{(b-a)^2}{12}.
\]
\end{stwierdzenie}
\begin{proof}
Gęstość $f(x)=(b-a)^{-1}\chi_{(a,b)}(x)$
jest borelowska, nieujemna i ma całkę
$(b-a)^{-1}\lambda_1((a,b))=1$, więc określa
miarę probabilistyczną.
Dla $x\leq a$ przekrój $(-\infty,x]\cap(a,b)$ jest pusty,
więc $F_X(x)=0$.
Dla $a<x<b$ całkujemy stałą gęstość od $a$ do $x$,
co daje $F_X(x)=(x-a)/(b-a)$.
Dla $x\geq b$ całka obejmuje cały nośnik i wynosi jeden.

Z faktu, że $\Prob(a<X<b)=1$, wynika ograniczoność $X$
prawie na pewno i skończoność wszystkich dodatnich momentów
bezwzględnych. Stosując \eqref{eq:lotus}, obliczamy:
\[
\E X=\frac1{b-a}\int_a^b x\,\dd\lambda_1(x)
=\frac{b^2-a^2}{2(b-a)}=\frac{a+b}2,
\]
\[
\E X^2=\frac1{b-a}\int_a^b x^2\,\dd\lambda_1(x)
=\frac{b^3-a^3}{3(b-a)}=\frac{a^2+ab+b^2}3.
\]
Korzystając ze wzoru na wariancję, otrzymujemy
\[
\Var(X)=\frac{a^2+ab+b^2}{3}-\frac{(a+b)^2}{4}
=\frac{4(a^2+ab+b^2)-3(a^2+2ab+b^2)}{12}
=\frac{(b-a)^2}{12}.
\]
\end{proof}

\begin{definicja}
Dla $\lambda>0$ rozkładem wykładniczym $\operatorname{Exp}(\lambda)$
z parametrem intensywności $\lambda$ nazywamy rozkład
o gęstości $f(x)=\lambda e^{-\lambda x}\chi_{[0,\infty)}(x)$
względem miary Lebesgue'a $\lambda_1$.
Jest to nieujemna funkcja borelowska, a jej całka wynosi
\[
\lim_{R\to\infty}\int_0^R\lambda e^{-\lambda x}\,
\dd\lambda_1(x)=\lim_{R\to\infty}(1-e^{-\lambda R})=1
\]
na mocy zbieżności monotonicznej, więc rzeczywiście definiuje
rozkład prawdopodobieństwa.
\end{definicja}

\begin{stwierdzenie}
Jeżeli $X\sim\operatorname{Exp}(\lambda)$, to dla $x\geq0$
$F_X(x)=1-e^{-\lambda x}$, a
$\E X=1/\lambda$ i $\Var(X)=1/\lambda^2$.
Dla $s,t\geq0$ zachodzi
$\Prob(X>s+t\mid X>s)=\Prob(X>t)$.
\end{stwierdzenie}
\begin{proof}
Dla $x<0$ mamy $F_X(x)=0$. Dla $x\geq0$ całkowanie
gęstości daje
$F_X(x)=\int_0^x\lambda e^{-\lambda t}\,\dd\lambda_1(t)
=1-e^{-\lambda x}$.
Zatem dla $x\geq0$ ogon ma postać
$\Prob(X>x)=e^{-\lambda x}$, a dla $x<0$
prawdopodobieństwo to wynosi jeden.

Aby obliczyć momenty, stosujemy \eqref{eq:lotus}
do funkcji nieujemnej $x\mapsto |x|^k$ i podstawiamy
$u=\lambda x$. Dla całkowitego $k\geq1$ mamy
\[
\E X^k=\E|X|^k
=\frac1{\lambda^k}\int_0^\infty u^ke^{-u}\,\dd\lambda_1(u).
\]
Całkowanie przez części na $[0,R]$ daje
\[
\int_0^R u^ke^{-u}\,\dd\lambda_1(u)
=-R^ke^{-R}+k\int_0^R u^{k-1}e^{-u}\,\dd\lambda_1(u).
\]
Ponieważ $R^ke^{-R}\to0$ i
$\int_0^\infty e^{-u}\,\dd\lambda_1(u)=1$,
indukcja i zbieżność monotoniczna dowodzą, że całka
jest skończona i równa $k!$. Stąd $\E X^k=k!/\lambda^k$.
Otrzymujemy $\E X=1/\lambda$ i
$\Var(X)=2/\lambda^2-1/\lambda^2=1/\lambda^2$.

Wreszcie $\{X>s+t\}\subseteq\{X>s\}$ oraz
$\Prob(X>s)=e^{-\lambda s}>0$, więc
\[
\Prob(X>s+t\mid X>s)
=\frac{\Prob(X>s+t)}{\Prob(X>s)}
=\frac{e^{-\lambda(s+t)}}{e^{-\lambda s}}
=e^{-\lambda t}=\Prob(X>t).
\]
Interpretacja: wiedza o dotychczasowym czasie oczekiwania
nie zmienia rozkładu dodatkowego czasu.
\end{proof}

\begin{uwaga}
Rozkład wykładniczy jest ciągłym odpowiednikiem rozkładu
geometrycznego ze względu na własność braku pamięci.
W procesach stochastycznych pojawi się jako rozkład czasów
oczekiwania w procesie Poissona.
\end{uwaga}

\begin{przyklad}[Generowanie zmiennej wykładniczej z jednostajnej]
Niech $U\sim\Unif(0,1)$ i $\lambda>0$. Definiujemy
$X=-\log(1-U)/\lambda$ na mierzalnym zdarzeniu $\{0<U<1\}$,
a na jego dopełnieniu przyjmujemy $X=0$.
Zmienna $X$ jest mierzalna, gdyż funkcja użyta w definicji
jest ciągła na $(0,1)$, a jej przedłużenie zerem jest borelowskie.
Ponieważ $\Prob(0<U<1)=1$, dla $x\geq0$ nierówność
$X\leq x$ jest równoważna $U\leq1-e^{-\lambda x}$
poza zbiorem zerowym. Zatem
$\Prob(X\leq x)=\Prob(U\leq1-e^{-\lambda x})=1-e^{-\lambda x}$.
Dla $x<0$ mamy $\Prob(X\leq x)=0$. Dystrybuanta jest
więc dystrybuantą $\operatorname{Exp}(\lambda)$, co
jednoznacznie wyznacza rozkład.
Jest to szczególny przypadek metody odwrotnej dystrybuanty,
opartej na kwantylu uogólnionym (rozdział~\ref{ch:random-variables}). Pozwala ona przejść od
jednostajnej liczby pseudolosowej do zadanej rodziny rozkładów.
\end{przyklad}

\section{Rozkład normalny}
Rozkład normalny wyróżnia jego postać analityczna oraz stabilność
względem dodawania niezależnych zmiennych normalnych. Służy też
jako rozkład graniczny dla odpowiednio unormowanych sum
wielu drobnych składników, co uzasadnimy później przy CTG.
Na razie nie będziemy motywować jego gęstości wyłącznie
rysunkiem dzwonu: wyprowadzimy normalizację przez całkę Gaussa
oraz sprawdzimy wzory na momenty i własności sum.

\begin{definicja}
Niech $\mu\in\R$ i $\sigma>0$. Rozkładem normalnym
$\Normal(\mu,\sigma^2)$ nazywamy rozkład o gęstości
\begin{equation}\label{eq:normal}
f(x)=\frac{1}{\sigma\sqrt{2\pi}}
\exp\!\left(-\frac{(x-\mu)^2}{2\sigma^2}\right),
\qquad x\in\R.
\end{equation}
Rozkład $\Normal(0,1)$ nazywamy standardowym rozkładem normalnym.
Jego dystrybuantę oznaczamy przez $\Phi$.
\end{definicja}

\begin{lemat}[Całka Gaussa]
Zachodzi $\int_\R e^{-x^2/2}\,\dd\lambda_1(x)=\sqrt{2\pi}$.
\end{lemat}
\begin{proof}
Niech $I=\int_\R e^{-x^2/2}\,\dd\lambda_1(x)\in(0,\infty]$.
Funkcja podcałkowa jest dodatnia i borelowska.
Miara $\lambda_1$ jest $\sigma$-skończona, a
$\lambda_2=\lambda_1\otimes\lambda_1$
na zbiorach borelowskich płaszczyzny
(por. \atlasref{ME-12::thm:lebesgue-products-12}{twierdzenie o produktach miar Lebesgue'a}).
Z \atlasref{ME-12::thm:tonelli-12}{twierdzenia Tonellego}
otrzymujemy równość $I^2$ z całką na $\R^2$, początkowo
w sensie rozszerzonym. Odwzorowanie
$(r,\varphi)\mapsto(r\cos\varphi,r\sin\varphi)$ jest
dyfeomorfizmem klasy $C^1$ z $(0,\infty)\times(0,2\pi)$
na płaszczyznę bez nieujemnej półosi poziomej;
jego macierz pochodnych ma postać
\[
\begin{pmatrix}
\cos\varphi&-r\sin\varphi\\
\sin\varphi&r\cos\varphi
\end{pmatrix},
\]
a moduł wyznacznika wynosi $r(\cos^2\varphi+\sin^2\varphi)=r$.
Pominięta półoś ma miarę Lebesgue'a $\lambda_2$ równą zero:
jest sumą przeliczalnie wielu odcinków zawartych
w prostokątach o dowolnie małej dodatniej wysokości.
Stosując \atlasref{ME-13::thm:change-variables-13}{twierdzenie o zamianie zmiennych}
do nieujemnej funkcji, otrzymujemy
\[
I^2=\int_{\R^2}e^{-(x^2+y^2)/2}\,\dd\lambda_2(x,y)
=\int_0^{2\pi}\!\int_0^\infty e^{-r^2/2}r\,
\dd\lambda_1(r)\,\dd\lambda_1(\varphi)=2\pi.
\]
W ostatnim kroku podstawiono $u=r^2/2$, co daje
$\int_0^\infty e^{-r^2/2}r\,\dd\lambda_1(r)
=\int_0^\infty e^{-u}\,\dd\lambda_1(u)=1$.
Ponieważ $I^2=2\pi$, całka $I$ jest skończona,
a z dodatniości wynika $I=\sqrt{2\pi}$.
W konsekwencji funkcja z \eqref{eq:normal} jest nieujemna,
borelowska i po podstawieniu $z=(x-\mu)/\sigma$
ma całkę $I/\sqrt{2\pi}=1$, więc definiuje rozkład.
\end{proof}

\begin{stwierdzenie}[Momenty i funkcja tworząca]
Jeżeli $X\sim\Normal(\mu,\sigma^2)$, to
$\E X=\mu$, $\Var(X)=\sigma^2$ oraz
$M_X(t)=\exp(\mu t+\sigma^2t^2/2)$ dla wszystkich $t\in\R$.
Ponadto $aX+b\sim\Normal(a\mu+b,a^2\sigma^2)$
dla $a\ne0$ i $b\in\R$; dla $a=0$ rozkład jest skupiony
w punkcie $b$.
\end{stwierdzenie}
\begin{proof}
Niech $Z=(X-\mu)/\sigma$. Z twierdzenia o
przekształceniach zmiennych (rozdział~\ref{ch:random-variables})
wynika, że $Z$ ma gęstość
$g(z)=\sigma f_X(\mu+\sigma z)=(2\pi)^{-1/2}e^{-z^2/2}$.
Korzystamy z dodatniości $\sigma$ oraz z faktu,
że odwrotność przekształcenia afinicznego ma pochodną $\sigma$.
Dla $t\in\R$, stosując \eqref{eq:lotus} do nieujemnej
funkcji $z\mapsto e^{t(\mu+\sigma z)}$, obliczamy
\[
M_X(t)=\frac{e^{\mu t}}{\sqrt{2\pi}}\int_\R
\exp\!\left(\sigma tz-\frac{z^2}2\right)
\,\dd\lambda_1(z).
\]
Tożsamość
$\sigma tz-z^2/2=-(z-\sigma t)^2/2+\sigma^2t^2/2$
pozwala wyjąć z całki czynnik
$e^{\sigma^2t^2/2}$. Pozostały czynnik wynosi jeden:
\[
\frac1{\sqrt{2\pi}}\int_\R e^{-(z-\sigma t)^2/2}\,
\dd\lambda_1(z)
=\frac1{\sqrt{2\pi}}\int_\R e^{-w^2/2}\,\dd\lambda_1(w)=1
\]
na mocy podstawienia $w=z-\sigma t$ i lematu o całce Gaussa.
Otrzymujemy więc
$M_X(t)=\exp(\mu t+\sigma^2t^2/2)$
dla wszystkich $t\in\R$.
Ponieważ funkcja ta jest skończona na całej prostej,
z twierdzenia o odzyskiwaniu momentów
(rozdział~\ref{ch:expectation}) wynika skończoność
wszystkich momentów bezwzględnych rzędów całkowitych oraz
$\E X=M_X'(0)$ i $\E X^2=M_X''(0)$.
Różniczkując otrzymany wzór, dostajemy
$M_X'(t)=(\mu+\sigma^2t)M_X(t)$ i
$M_X''(t)=[\sigma^2+(\mu+\sigma^2t)^2]M_X(t)$.
Ponieważ $M_X(0)=1$, mamy $\E X=\mu$ i
$\E X^2=\mu^2+\sigma^2$, zatem $\Var(X)=\sigma^2$.
Dla rzeczywistego $a\ne0$ transformacja gęstości
przez funkcję $x\mapsto ax+b$ daje
\[
f_{aX+b}(y)=\frac1{|a|}f_X\!\left(\frac{y-b}{a}\right)
=\frac1{|a|\sigma\sqrt{2\pi}}
\exp\!\left(-\frac{(y-a\mu-b)^2}{2a^2\sigma^2}\right).
\]
Jest to gęstość $\Normal(a\mu+b,a^2\sigma^2)$;
gdy $a=0$, zmienna $aX+b$ jest stale równa $b$.
\end{proof}

\begin{stwierdzenie}[Suma niezależnych zmiennych normalnych]
Jeżeli $X\sim\Normal(\mu_1,\sigma_1^2)$ i
$Y\sim\Normal(\mu_2,\sigma_2^2)$ są niezależne na
wspólnej przestrzeni probabilistycznej, gdzie
$\mu_1,\mu_2\in\R$ i $\sigma_1,\sigma_2>0$, to
$X+Y\sim\Normal(\mu_1+\mu_2,\sigma_1^2+\sigma_2^2)$.
\end{stwierdzenie}
\begin{proof}
Z niezależności i wzoru na splot gęstości
(rozdział~\ref{ch:random-variables}) wynika, że następująca
funkcja jest wersją gęstości sumy dla $s\in\R$:
\[
f_{X+Y}(s)=\frac1{2\pi\sigma_1\sigma_2}
\int_\R\exp\!\left(-\frac{(x-\mu_1)^2}{2\sigma_1^2}
-\frac{(s-x-\mu_2)^2}{2\sigma_2^2}\right)\,\dd\lambda_1(x).
\]
Oznaczmy $v=\sigma_1^2+\sigma_2^2$ oraz
$c(s)=(\mu_1\sigma_2^2+(s-\mu_2)\sigma_1^2)/v$.
Uzupełnienie kwadratu daje
\[
\frac{(x-\mu_1)^2}{\sigma_1^2}
+\frac{(s-x-\mu_2)^2}{\sigma_2^2}
=\frac{v}{\sigma_1^2\sigma_2^2}(x-c(s))^2
+\frac{(s-\mu_1-\mu_2)^2}{v}.
\]
Istotnie, pisząc $A=\sigma_1^2$, $B=\sigma_2^2$,
$m=\mu_1$ i $n=s-\mu_2$, otrzymujemy
\[
\frac{(x-m)^2}{A}+\frac{(x-n)^2}{B}
=\frac{(A+B)x^2-2(Bm+An)x+Bm^2+An^2}{AB}.
\]
Po wydzieleniu kwadratu z $c=(Bm+An)/(A+B)$
pozostaje wyraz stały
\[
\frac{(A+B)(Bm^2+An^2)-(Bm+An)^2}{AB(A+B)}
=\frac{(m-n)^2}{A+B},
\]
co daje powyższą tożsamość.
Po wyłączeniu czynnika niezależnego od $x$ stosujemy
podstawienie $u=\sqrt v(x-c(s))/(\sigma_1\sigma_2)$.
Całka po $\R$ z pozostałego czynnika wynosi
$\sigma_1\sigma_2\sqrt{2\pi}/\sqrt v$ na mocy lematu
o całce Gaussa. W konsekwencji otrzymujemy
\[
f_{X+Y}(s)=\frac{1}{\sqrt{2\pi v}}
\exp\!\left(-\frac{(s-\mu_1-\mu_2)^2}{2v}\right),
\]
czyli żądaną gęstość.
\end{proof}

\begin{przyklad}[Standaryzacja]
Jeżeli $X\sim\Normal(\mu,\sigma^2)$, to
$Z:=(X-\mu)/\sigma\sim\Normal(0,1)$.
W konsekwencji dla $a<b$
$\{a<X\leq b\}=\{(a-\mu)/\sigma<Z\leq(b-\mu)/\sigma\}$,
gdyż $\sigma>0$. Z definicji dystrybuanty
\[
\Prob(a<X\leq b)=\Phi((b-\mu)/\sigma)-\Phi((a-\mu)/\sigma).
\]
Funkcja $\Phi(x)=(2\pi)^{-1/2}\int_{-\infty}^x
e^{-u^2/2}\,\dd\lambda_1(u)$ jest funkcją pierwotną
standardowej gęstości normalnej na mocy twierdzenia
podstawowego rachunku całkowego dla funkcji ciągłych
(po odjęciu stałej $\Phi(0)$ całkujemy od $0$ do $x$). Jej wartości wyznacza
się numerycznie lub za pomocą funkcji błędu.
\end{przyklad}

\section{Funkcja gamma i funkcja beta}
Próby normalizowania funkcji takich jak
$x^{a-1}e^{-x}$ albo $x^{a-1}(1-x)^{b-1}$ prowadzą
do całek, które nie zawsze są elementarne. Funkcje gamma i beta
zapisują te całki w sposób systematyczny. Nie są jedynie
technicznym dodatkiem: ich tożsamości pozwalają wyprowadzać
wzory na momenty i zależności między rozkładami gamma,
chi-kwadrat, beta i t-Studenta.

\begin{definicja}
Dla $a,b>0$ definiujemy funkcję gamma i funkcję beta Eulera wzorami
\[
\Gamma(a):=\int_0^\infty t^{a-1}e^{-t}\,\dd\lambda_1(t),
\qquad
\mathrm B(a,b):=\int_0^1 t^{a-1}(1-t)^{b-1}\,\dd\lambda_1(t).
\]
Całki rozumiemy jako całki Lebesgue'a po odpowiednio
$(0,\infty)$ i $(0,1)$, bez potrzeby określania wartości
funkcji podcałkowych w punktach końcowych.
Obie całki są skończone i dodatnie dla podanych parametrów,
co uzasadnimy w następnym dowodzie.
\end{definicja}

\begin{lemat}[Podstawowe tożsamości]
Dla $a,b>0$ zachodzą
$\Gamma(a+1)=a\Gamma(a)$ oraz
$\mathrm B(a,b)=\Gamma(a)\Gamma(b)/\Gamma(a+b)$.
Ponadto $\Gamma(1)=1$, $\Gamma(1/2)=\sqrt\pi$
i $\Gamma(n+1)=n!$ dla $n\in\N\cup\{0\}$.
\end{lemat}
\begin{proof}
Najpierw uzasadnijmy zbieżność całki definiującej
$\Gamma(a)$. Na $(0,1)$ mamy $0<t^{a-1}e^{-t}\leq t^{a-1}$,
a $\int_0^1t^{a-1}\,\dd\lambda_1(t)=1/a<\infty$.
Na $[1,\infty)$ funkcja $t^{a-1}e^{-t/2}$ jest
ciągła i ograniczona: jej granica w nieskończoności
jest zerowa, gdyż $(a-1)\log t-t/2\to-\infty$.
Zatem $t^{a-1}e^{-t}\leq C_a e^{-t/2}$ dla pewnej
skończonej stałej $C_a$; prawa strona jest całkowalna.
W przypadku funkcji beta na $(0,1/2]$ czynnik
$(1-t)^{b-1}$ jest ograniczony przez $\max(1,2^{1-b})$,
a $\int_0^{1/2}t^{a-1}\,\dd\lambda_1(t)=2^{-a}/a$.
Na $[1/2,1)$ czynnik $t^{a-1}$ jest ograniczony przez
$\max(1,2^{1-a})$, a całka z $(1-t)^{b-1}$
wynosi $2^{-b}/b$. Otrzymujemy skończoność obu całek.
Ich ścisła dodatniość wynika z dodatniości funkcji
podcałkowych na przedziałach dodatniej miary.

Całkując przez części na $[\varepsilon,R]$,
gdzie $0<\varepsilon<R$, otrzymujemy
\[
\int_\varepsilon^R t^ae^{-t}\,\dd\lambda_1(t)
=-R^ae^{-R}+\varepsilon^ae^{-\varepsilon}
+a\int_\varepsilon^R t^{a-1}e^{-t}\,\dd\lambda_1(t).
\]
Przy $\varepsilon\downarrow0$ i $R\uparrow\infty$
wyrazy brzegowe dążą do zera (ponieważ $a>0$
i wykładnicza dominuje potęgę w nieskończoności).
Zbieżność monotoniczna dla nieujemnych funkcji
podcałkowych pozwala przejść do granicy. Otrzymujemy
\[
\Gamma(a+1)=\int_0^\infty t^ae^{-t}\,\dd\lambda_1(t)
=a\int_0^\infty t^{a-1}e^{-t}\,\dd\lambda_1(t)
=a\Gamma(a).
\]
Ponieważ $\Gamma(1)=\int_0^\infty e^{-t}\,\dd\lambda_1(t)=1$,
indukcja daje $\Gamma(n+1)=n!$.
Podstawienie $t=u^2$ jest zmianą zmiennych na
$(0,\infty)$ o jakobianie $2u$, więc
$\Gamma(1/2)=2\int_0^\infty e^{-u^2}\,\dd\lambda_1(u)$.
Z symetrii i podstawienia $v=\sqrt2\,u$ oraz
wcześniejszego lematu o całce Gaussa wynika
\[
2\int_0^\infty e^{-u^2}\,\dd\lambda_1(u)
=\frac1{\sqrt2}\int_\R e^{-v^2/2}\,\dd\lambda_1(v)=\sqrt\pi.
\]

Aby wyprowadzić relację z funkcją beta, stosujemy
\atlasref{ME-12::thm:tonelli-12}{twierdzenie Tonellego}
do nieujemnej borelowskiej funkcji na $(0,\infty)^2$
względem miary produktowej Lebesgue'a:
\[
\Gamma(a)\Gamma(b)
=\int_{(0,\infty)^2}u^{a-1}v^{b-1}e^{-u-v}
\,\dd\lambda_2(u,v).
\]
Wprowadzamy $r=u+v>0$ i $w=u/(u+v)\in(0,1)$.
Odwrotnie, $u=rw$, $v=r(1-w)$.
Oba odwzorowania są klasy $C^1$ na odpowiednich
otwartych obszarach, więc jest to dyfeomorfizm
$(0,\infty)\times(0,1)$ na $(0,\infty)^2$.
Wyznacznik macierzy pochodnych odwzorowania
$(r,w)\mapsto(u,v)$ wynosi
\[
\det\begin{pmatrix}w&r\\1-w&-r\end{pmatrix}=-r.
\]
Z \atlasref{ME-13::thm:change-variables-13}{twierdzenia o zamianie zmiennych}
funkcja podcałkowa pomnożona przez moduł jakobianu
przyjmuje postać
$r^{a+b-1}e^{-r}w^{a-1}(1-w)^{b-1}$.
Oba czynniki są nieujemne i borelowskie.
Ponowne zastosowanie twierdzenia Tonellego pozwala
rozdzielić całkę:
\[
\Gamma(a)\Gamma(b)
=\left(\int_0^\infty r^{a+b-1}e^{-r}\,\dd\lambda_1(r)\right)
\left(\int_0^1w^{a-1}(1-w)^{b-1}\,\dd\lambda_1(w)\right)
=\Gamma(a+b)\mathrm B(a,b).
\]
Dzielenie przez dodatnie $\Gamma(a+b)$ kończy dowód.
\end{proof}

\section{Rozkład gamma}
Suma skończonej liczby niezależnych zmiennych o tym samym
rozkładzie wykładniczym prowadzi do rodziny gamma,
co wynika z poniższej własności dodawania. Parametr kształtu opisuje m.in. liczbę składników
w przypadkach całkowitych, a parametr intensywności reguluje
skalę czasową. Funkcja gamma spełnia $\Gamma(n+1)=n!$ dla
całkowitego $n\geq0$, więc funkcja $x\mapsto\Gamma(x+1)$
przedłuża silnię na rzeczywiste argumenty $x\geq0$.

\begin{definicja}
Niech $\alpha,\lambda>0$. Rozkładem gamma $\Gam(\alpha,\lambda)$
o parametrze kształtu $\alpha$ i \emph{intensywności} $\lambda$
nazywamy rozkład o gęstości
\begin{equation}\label{eq:gamma}
f(x)=
\begin{cases}
\displaystyle\frac{\lambda^\alpha}{\Gamma(\alpha)}
 x^{\alpha-1}e^{-\lambda x},&x>0,\\
0,&x\leq0.
\end{cases}
\end{equation}
\end{definicja}

\begin{stwierdzenie}
Dla $X\sim\Gam(\alpha,\lambda)$ zachodzi
\[
\E X=\frac\alpha\lambda,\qquad
\Var(X)=\frac\alpha{\lambda^2},\qquad
M_X(t)=\left(\frac\lambda{\lambda-t}\right)^\alpha\quad(t<\lambda).
\]
Dla $t\geq\lambda$ zachodzi $M_X(t)=\infty$.
Jeżeli $X\sim\Gam(\alpha,\lambda)$ i
$Y\sim\Gam(\beta,\lambda)$, gdzie $\beta>0$, są niezależne
na wspólnej przestrzeni probabilistycznej, to
$X+Y\sim\Gam(\alpha+\beta,\lambda)$.
\end{stwierdzenie}
\begin{proof}
Najpierw sprawdźmy normalizację gęstości, która
jest borelowska i nieujemna (również przy $\alpha<1$,
ponieważ potęga jest obliczana tylko dla $x>0$).
Dla $\alpha,\lambda>0$, po podstawieniu $u=\lambda x$,
\[
\int_0^\infty\frac{\lambda^\alpha}{\Gamma(\alpha)}
 x^{\alpha-1}e^{-\lambda x}\,\dd\lambda_1(x)
=\frac1{\Gamma(\alpha)}\int_0^\infty
u^{\alpha-1}e^{-u}\,\dd\lambda_1(u)=1.
\]
Dla rzeczywistego $r\geq0$ ze wzoru \eqref{eq:lotus}
oraz tego samego podstawienia otrzymujemy
\[
\E X^r
=\frac{\lambda^\alpha}{\Gamma(\alpha)}
\int_0^\infty x^{\alpha+r-1}e^{-\lambda x}\,\dd\lambda_1(x)
=\frac{\Gamma(\alpha+r)}{\lambda^r\Gamma(\alpha)}.
\]
Potęgę $X^r$ dla niecałkowitego $r$ rozumiemy jako
funkcję określoną przez $x^r$ dla $x>0$ i przedłużoną
mierzalnie przez zero na zdarzeniu $\{X\leq0\}$
o prawdopodobieństwie zero. Dla $r=0$ przyjmujemy $X^0=1$.
Całka jest skończona, gdyż $\alpha+r>0$.
Korzystając z rekurencji funkcji gamma dla $r=1,2$,
mamy $\E X=\alpha/\lambda$ i
$\E X^2=\alpha(\alpha+1)/\lambda^2$.
Różnica drugiego momentu i kwadratu średniej daje
$\Var(X)=\alpha/\lambda^2$.

Dla $t<\lambda$ obliczamy
\[
M_X(t)=\frac{\lambda^\alpha}{\Gamma(\alpha)}
\int_0^\infty x^{\alpha-1}e^{-(\lambda-t)x}
\,\dd\lambda_1(x)
=\left(\frac{\lambda}{\lambda-t}\right)^\alpha,
\]
gdzie użyliśmy kolejno zmiany $u=(\lambda-t)x$
i definicji funkcji gamma (argument $\alpha$ jest dodatni).
Dla $t\geq\lambda$ mamy $e^{-(\lambda-t)x}\geq1$ na
$[1,\infty)$, a $\int_1^\infty x^{\alpha-1}\,\dd\lambda_1(x)
=\infty$ dla $\alpha>0$. Wobec tego funkcja tworząca
momenty ma dziedzinę skończoności dokładnie $(-\infty,\lambda)$.

Dla sumy niezależnych zmiennych gamma z tą samą
intensywnością korzystamy ze wzoru na splot
(rozdział~\ref{ch:random-variables}).
Dla $s>0$ obie gęstości są dodatnie jedynie wtedy,
gdy jednocześnie $x>0$ i $s-x>0$, czyli $x\in(0,s)$.
Dla $s\leq0$ taki $x$ nie istnieje i splot jest zerowy.
Dla $s>0$ otrzymujemy:
\[
f_{X+Y}(s)=\frac{\lambda^{\alpha+\beta}e^{-\lambda s}}
{\Gamma(\alpha)\Gamma(\beta)}
\int_0^s x^{\alpha-1}(s-x)^{\beta-1}\,\dd\lambda_1(x),
\quad s>0.
\]
Podstawiamy $x=sw$ dla $0<w<1$; moduł jakobianu
wynosi $s$. Stąd
\[
\int_0^s x^{\alpha-1}(s-x)^{\beta-1}\,\dd\lambda_1(x)
=s^{\alpha+\beta-1}\mathrm B(\alpha,\beta).
\]
Całka jest skończona, gdyż oba parametry są dodatnie.
Po zastosowaniu tożsamości beta--gamma otrzymujemy
$f_{X+Y}(s)=\lambda^{\alpha+\beta}s^{\alpha+\beta-1}
e^{-\lambda s}/\Gamma(\alpha+\beta)$ dla $s>0$,
a dla $s\leq0$ gęstość jest zerowa. Jest to dokładnie
gęstość $\Gam(\alpha+\beta,\lambda)$.
Wspólna intensywność jest istotna: dla różnych
intensywności ten wzór na sumę nie obowiązuje.
Na przykład dla dwóch niezależnych zmiennych wykładniczych
o intensywnościach $a\ne b$, $a,b>0$, średnia sumy
wynosi $1/a+1/b$, a wariancja $1/a^2+1/b^2$.
Używamy oznaczeń $a,b$, aby nie mylić parametrów
z oznaczeniami miar Lebesgue'a.
Gdyby suma miała rozkład gamma o kształcie $2$,
jej wariancja musiałaby być równa połowie kwadratu średniej.
Różnica między rzeczywistą wariancją i tą wartością wynosi
$\tfrac12(1/a-1/b)^2>0$.
\end{proof}

\begin{uwaga}
W literaturze używa się również parametru \emph{skali}
$\beta=1/\lambda$, zapisując gęstość jako
$x^{\alpha-1}e^{-x/\beta}/(\Gamma(\alpha)\beta^\alpha)$.
W tej książce drugi parametr rozkładu gamma jest zawsze
\emph{intensywnością}, a nie skalą.
Rozkład $\operatorname{Exp}(\lambda)$ jest równy $\Gam(1,\lambda)$,
ponieważ dla $x>0$ ich gęstości są identyczne
($\Gamma(1)=1$), a ewentualna różnica w punkcie $0$
nie wpływa na rozkład.
W konsekwencji, jeżeli $X_1,\ldots,X_n$ są niezależne
i mają rozkład $\operatorname{Exp}(\lambda)$, gdzie
$n\geq1$ jest całkowite, to $\sum_{j=1}^n X_j\sim\Gam(n,\lambda)$.
Istotnie, stosujemy indukcyjnie stwierdzenie o sumie
zmiennych gamma, korzystając z niezależności sumy
poprzednich składników od następnego składnika.
\end{uwaga}

\section{Rozkład chi-kwadrat}
Suma kwadratów niezależnych standardowych zmiennych normalnych
jest naturalną miarą kwadratowego odchylenia wektora losowego
od początku układu. Jej rozkład należy do rodziny gamma i
nazywa się chi-kwadrat. Dla całkowitej liczby stopni swobody
ma ona bezpośredni sens geometryczny; definicję przez gamma
można natomiast rozszerzyć na dowolne dodatnie $\nu$.

\begin{definicja}
Niech $\nu>0$. Rozkładem chi-kwadrat z $\nu$ stopniami swobody,
oznaczanym $\chi^2_\nu$, nazywamy rozkład
$\Gam(\nu/2,1/2)$, a więc rozkład o gęstości
\begin{equation}\label{eq:chisq}
f_\nu(x)=
\begin{cases}
\displaystyle\frac{1}{2^{\nu/2}\Gamma(\nu/2)}
 x^{\nu/2-1}e^{-x/2},&x>0,\\
0,&x\leq0.
\end{cases}
\end{equation}
Dla całkowitego $\nu\geq1$ parametr $\nu$ oznacza liczbę
niezależnych standardowych zmiennych normalnych
w poniższej reprezentacji.
\end{definicja}

\begin{twierdzenie}[Suma kwadratów zmiennych normalnych]
Niech $\nu\geq1$ będzie liczbą całkowitą i niech
$Z_1,\ldots,Z_\nu$ będą niezależnymi zmiennymi losowymi
na wspólnej przestrzeni probabilistycznej, o rozkładzie
$\Normal(0,1)$. Wówczas
$U:=\sum_{j=1}^\nu Z_j^2\sim\chi^2_\nu$.
Dla dowolnego rzeczywistego $\nu>0$, jeżeli $U\sim\chi^2_\nu$,
to $\E U=\nu$ oraz $\Var(U)=2\nu$.
\end{twierdzenie}
\begin{proof}
Wystarczy najpierw obliczyć rozkład $Z^2$ dla jednej
standardowej zmiennej normalnej $Z$.
Niech $\varphi(z)=(2\pi)^{-1/2}e^{-z^2/2}$ oznacza
gęstość $Z$. Funkcja $g(z)=z^2$ jest dyfeomorfizmem
klasy $C^1$ z każdego z przedziałów $(-\infty,0)$
i $(0,\infty)$ na $(0,\infty)$.
Pominięty punkt $0$ ma prawdopodobieństwo zero,
ponieważ $Z$ ma gęstość. Dla $u>0$ przeciwobrazy to
$z_\pm=\pm\sqrt u$, a $|g'(z_\pm)|=2\sqrt u>0$.
Ze wzoru na transformację o dwóch gałęziach
(rozdział~\ref{ch:random-variables})
\[
f_{Z^2}(u)=\frac{\varphi(\sqrt u)}{2\sqrt u}
+\frac{\varphi(-\sqrt u)}{2\sqrt u}
=\frac{e^{-u/2}}{\sqrt{2\pi u}}
\qquad(u>0).
\]
Ponieważ $\Gamma(1/2)=\sqrt\pi$, wzór ten jest równy
$u^{-1/2}e^{-u/2}/(2^{1/2}\Gamma(1/2))$.
Dla $u\leq0$ przyjmujemy gęstość równą zeru,
co jest dopuszczalne, ponieważ $Z^2>0$ prawie na pewno.
Jest to gęstość $\Gam(1/2,1/2)$.

Każda zmienna $Z_j^2$ ma więc taki sam rozkład gamma.
Dla dowolnych borelowskich $B_j\subseteq\R$
zdarzenie $\{Z_j^2\in B_j\}$ jest zdarzeniem
$\{Z_j\in g^{-1}(B_j)\}$, gdzie przeciwobraz jest
borelowski. Z niezależności $Z_j$ wynika zatem
niezależność $Z_j^2$. Kolejne sumy tych zmiennych
są niezależne od następnego składnika, ponieważ
rozkład łączny jest miarą produktową.
Indukcyjne zastosowanie własności dodawania niezależnych
rozkładów gamma o wspólnej intensywności daje
$\sum_{j=1}^\nu Z_j^2\sim\Gam(\nu/2,1/2)=\chi^2_\nu$.
Niezależnie od tego, czy $\nu$ jest całkowite,
przez definicję $\chi^2_\nu$ jako rozkładu gamma
obliczamy $\E U=(\nu/2)/(1/2)=\nu$ oraz
$\Var(U)=(\nu/2)/(1/2)^2=2\nu$.
\end{proof}

\begin{przyklad}[Trzy współrzędne normalne]
Jeżeli $Z_1,Z_2,Z_3$ są niezależne i standardowo normalne,
to $R^2:=Z_1^2+Z_2^2+Z_3^2\sim\chi^2_3$.
Średnia $R^2$ wynosi $3$, a wariancja $6$.
Nie należy mylić rozkładu $R^2$ z rozkładem $R$:
po spierwiastkowaniu powstaje inny rozkład, określany
jako rozkład chi.
\end{przyklad}

\section{Rozkład beta}
W przeciwieństwie do rozkładu gamma, skupionego na dodatniej
półprostej, rozkład beta jest skupiony na przedziale $(0,1)$.
Dzięki dwóm parametrom jego gęstość może być symetryczna,
mieć maksimum wewnątrz przedziału lub rosnąć przy jednym
bądź obu końcach.
Jest więc użyteczny do modelowania proporcji i udziałów.
Jego normalizacja pochodzi wprost z funkcji beta Eulera.

\begin{definicja}
Niech $a,b>0$. Rozkładem beta $\Betad(a,b)$ na $(0,1)$
nazywamy rozkład o gęstości
\begin{equation}\label{eq:beta}
f(x)=
\begin{cases}
\displaystyle\frac{x^{a-1}(1-x)^{b-1}}{\mathrm B(a,b)},&0<x<1,\\
0,&x\notin(0,1).
\end{cases}
\end{equation}
\end{definicja}

\begin{stwierdzenie}
Jeżeli $X\sim\Betad(a,b)$, to
\[
\E X=\frac a{a+b},\qquad
\Var(X)=\frac{ab}{(a+b)^2(a+b+1)}.
\]
Rozkład $\Betad(1,1)$ jest jednostajny na $(0,1)$.
\end{stwierdzenie}
\begin{proof}
Najpierw sprawdzamy normalizację: gęstość jest
borelowska i nieujemna, a jej całka po $\R$ równa
$\mathrm B(a,b)/\mathrm B(a,b)=1$, gdyż $a,b>0$.
Zatem $\Prob(0<X<1)=1$, co zapewnia skończoność
wszystkich momentów bezwzględnych dodatnich rzędów.
Dla rzeczywistego $r\geq0$ stosujemy \eqref{eq:lotus}
do funkcji $x\mapsto x^r$ na $(0,1)$, przedłużonej
mierzalnie przez zero poza tym przedziałem.
Dla $r=0$ przyjmujemy $X^0=1$ prawie na pewno.
Z definicji funkcji beta i tożsamości beta--gamma wynika
\[
\E X^r=\frac{\mathrm B(a+r,b)}{\mathrm B(a,b)}
=\frac{\Gamma(a+r)\Gamma(a+b)}
{\Gamma(a)\Gamma(a+b+r)}.
\]
Dla $r=1$ rekurencja $\Gamma(s+1)=s\Gamma(s)$
prowadzi do $\E X=a/(a+b)$. Dla $r=2$
otrzymujemy
$\E X^2=a(a+1)/[(a+b)(a+b+1)]$.
Stąd
\[
\Var(X)=\frac{a(a+1)}{(a+b)(a+b+1)}
-\frac{a^2}{(a+b)^2}
=\frac{ab}{(a+b)^2(a+b+1)}.
\]
Wreszcie $\mathrm B(1,1)=\int_0^1 1\,\dd\lambda_1(x)=1$,
więc dla $a=b=1$ gęstość jest stała i równa jeden
na $(0,1)$, tak jak w rozkładzie jednostajnym.
\end{proof}

\section{Rozkład t-Studenta}
Jeżeli zmienną normalną dzielimy przez niezależne losowe
odchylenie związane z rozkładem chi-kwadrat, mianownik może
przyjmować małe wartości. Powstają wówczas obserwacje znacznie
bardziej odległe od zera niż w rozkładzie normalnym. To źródło
cięższych ogonów rozkładu t-Studenta. Liczba stopni swobody
steruje częstością takich odchyleń i decyduje, które momenty
w ogóle istnieją.

\begin{definicja}
Niech $\nu>0$, a $Z\sim\Normal(0,1)$ oraz
$U\sim\chi^2_\nu$ będą niezależnymi zmiennymi na wspólnej
przestrzeni probabilistycznej (takie zmienne można
skonstruować na iloczynie przestrzeni z tymi rozkładami).
Ponieważ $\Prob(U>0)=1$, definiujemy
\[
T=\begin{cases}
Z/\sqrt{U/\nu},&U>0,\\
0,&U\leq0.
\end{cases}
\]
Jest to zmienna mierzalna: funkcja dwóch zmiennych
użyta w definicji jest ciągła na $\R\times(0,\infty)$
i przedłużona zerem na dopełnieniu tego zbioru.
Jej rozkład zależy wyłącznie od $\nu$, gdyż rozkład
łączny $(Z,U)$ jest iloczynem ustalonych rozkładów,
a rozkład $T$ jest jego obrazem.
Nazywamy go rozkładem t-Studenta z $\nu$ stopniami swobody
i zapisujemy $T\sim t_\nu$.
\end{definicja}

\begin{twierdzenie}[Gęstość rozkładu Studenta]
Dla każdego $\nu>0$ zmienna $T\sim t_\nu$ ma gęstość
\begin{equation}\label{eq:student}
f_\nu(t)=
\frac{\Gamma((\nu+1)/2)}{\sqrt{\nu\pi}\,\Gamma(\nu/2)}
\left(1+\frac{t^2}{\nu}\right)^{-(\nu+1)/2},
\qquad t\in\R.
\end{equation}
Gęstość jest parzysta. Dla $\nu>1$ wartość oczekiwana
jest równa zero, a dla $\nu>2$ wariancja wynosi $\nu/(\nu-2)$.
Dla dowolnego rzeczywistego $r>0$ moment bezwzględny
$\E|T|^r$ jest skończony wtedy i tylko wtedy, gdy $r<\nu$.
Dla $0<\nu\leq1$ obie części całki definiującej wartość
oczekiwaną są nieskończone, więc wartość oczekiwana nie
istnieje. Dla $0<\nu\leq2$ drugi moment jest nieskończony.
\end{twierdzenie}
\begin{proof}
Wystarczy rozpatrzyć zmienną $T$ skonstruowaną z niezależnych
$Z\sim\Normal(0,1)$ i $U\sim\chi_\nu^2$ w definicji,
gdyż gęstość i momenty zależą wyłącznie od rozkładu.
Niech $\varphi(z)=(2\pi)^{-1/2}e^{-z^2/2}$ oznacza
gęstość standardowego rozkładu normalnego.
Z niezależności i wzoru na gęstość rozkładu łącznego
(rozdział~\ref{ch:random-variables}) para $(Z,U)$
ma gęstość $\varphi(z)f_U(u)$, gdzie dla $u>0$
$f_U(u)=u^{\nu/2-1}e^{-u/2}/(2^{\nu/2}\Gamma(\nu/2))$,
a dla $u\leq0$ gęstość jest zerowa.
Odwzorowanie $(z,u)\mapsto(z\sqrt{\nu/u},u)$
jest dyfeomorfizmem klasy $C^1$ obszaru
$\R\times(0,\infty)$ na siebie, z odwrotnością
$(t,u)\mapsto(t\sqrt{u/\nu},u)$.
Macierz pochodnych odwrotności jest równa
\[
\begin{pmatrix}
\sqrt{u/\nu}&t/(2\sqrt{\nu u})\\
0&1
\end{pmatrix},
\]
więc moduł jej wyznacznika wynosi $\sqrt{u/\nu}$.
Rozkład $(Z,U)$ jest skupiony w tym obszarze,
zatem \atlasref{ME-13::thm:change-variables-13}{twierdzenie o zamianie zmiennych}
daje gęstość pary $(T,U)$ równą
$\varphi(t\sqrt{u/\nu})f_U(u)\sqrt{u/\nu}$ dla $u>0$
i zeru poza tym obszarem.
Stosując \atlasref{ME-12::thm:tonelli-12}{twierdzenie Tonellego}
do nieujemnej gęstości rozkładu łącznego, uzyskujemy
gęstość brzegową $T$ przez całkowanie po $u$.
Jej całka jest równa jeden, ponieważ całka gęstości
rozkładu łącznego $(T,U)$ wynosi jeden.
Po zebraniu potęg otrzymujemy
\[
f_T(t)=\frac1{\sqrt{2\pi\nu}\,2^{\nu/2}\Gamma(\nu/2)}
\int_0^\infty u^{(\nu+1)/2-1}
\exp\!\left(-\frac{1+t^2/\nu}{2}u\right)
\,\dd\lambda_1(u).
\]
Oznaczmy $a=(\nu+1)/2$ i $c=(1+t^2/\nu)/2>0$.
Podstawienie $v=cu$ daje
$\int_0^\infty u^{a-1}e^{-cu}\,\dd\lambda_1(u)
=\Gamma(a)c^{-a}$.
Po wstawieniu i uproszczeniu potęg dwójki
otrzymujemy
\[
f_T(t)=\frac{\Gamma((\nu+1)/2)}
{\sqrt{\nu\pi}\,\Gamma(\nu/2)}
\left(1+\frac{t^2}{\nu}\right)^{-(\nu+1)/2}.
\]
Gęstość jest ciągła, dodatnia i parzysta, ale symetria
sama w sobie nie gwarantuje istnienia wartości oczekiwanej.
Na $[-1,1]$ funkcja $|t|^rf_T(t)$ dla $r\geq0$
jest ograniczona (dla $r=0$ przyjmujemy $|t|^0=1$).
Dla $|t|\geq1$ zachodzi
\[
\frac{t^2}{\nu}\leq1+\frac{t^2}{\nu}
\leq(1+1/\nu)t^2.
\]
Podnosząc strony do potęgi $-(\nu+1)/2$ i mnożąc
przez dodatni współczynnik z \eqref{eq:student},
otrzymujemy stałe $c_1,c_2>0$ zależne wyłącznie od $\nu$
takie, że $c_1|t|^{-\nu-1}\leq f_T(t)
\leq c_2|t|^{-\nu-1}$. Zatem dla $r\geq0$
całka $\int_{|t|\geq1}|t|^rf_T(t)\,\dd\lambda_1(t)$
zachowuje się pod względem zbieżności jak
$\int_1^\infty t^{r-\nu-1}\,\dd\lambda_1(t)$,
a ta jest skończona dokładnie dla $r<\nu$.
W konsekwencji $\E|T|<\infty$ wtedy i tylko wtedy,
gdy $\nu>1$; w tym przypadku $\E T=0$ przez
nieparzystość całkowalnej funkcji $t\mapsto tf_T(t)$.
Dla $0<\nu\leq1$ całki
$\int_0^\infty tf_T(t)\,\dd\lambda_1(t)$ oraz
$\int_{-\infty}^0(-t)f_T(t)\,\dd\lambda_1(t)$
są równe z parzystości gęstości, a ich suma jest
nieskończona, więc obie są nieskończone.

Dla $\nu>2$ obliczamy odwrotny moment $U$, przyjmując
$U^{-1}=0$ na $\{U\leq0\}$. Podstawienie $v=u/2$
i rekurencja funkcji gamma dają
\[
\E(U^{-1})=
\frac{1}{2^{\nu/2}\Gamma(\nu/2)}
\int_0^\infty u^{\nu/2-2}e^{-u/2}\,\dd\lambda_1(u)
=\frac{\Gamma(\nu/2-1)}{2\Gamma(\nu/2)}
=\frac1{\nu-2}.
\]
Argument $\nu/2-1$ funkcji gamma jest dodatni,
więc całka jest skończona. Ponieważ $Z,U$ są niezależne,
rozkład łączny jest produktowy; Tonelli zastosowany
do nieujemnej funkcji $(z,u)\mapsto z^2/u$ dla $u>0$
(zerowej dla $u\leq0$) daje
\[
\E T^2=\nu\E Z^2\,\E(U^{-1})
=\frac{\nu}{\nu-2},
\]
gdzie $\E Z^2=1$ z własności rozkładu normalnego.
Stąd $\Var(T)=\nu/(\nu-2)$.
Dla $1<\nu\leq2$ średnia wynosi zero, lecz drugi moment
jest nieskończony. Nie definiujemy wtedy wariancji
w sensie definicji z rozdziału~\ref{ch:expectation},
która wymaga skończonego drugiego momentu.
Jeżeli dopuścimy wariancję rozszerzoną
$\E[(T-\E T)^2]$, wynosi ona $+\infty$.
Dla $0<\nu\leq1$ wariancja nie jest określona,
gdyż nie istnieje średnia.
\end{proof}

\begin{uwaga}
Parametr $\nu$ w definicji rozkładu Studenta może być dowolną
dodatnią liczbą rzeczywistą. W klasycznych zadaniach o estymacji
wariancji na podstawie próby normalnej występują zwykle całkowite
liczby stopni swobody, często $n-1$.
Dla $\nu\to\infty$ rozkłady $t_\nu$ zbiegają według
rozkładu do $\Normal(0,1)$. Dla każdego $\nu>0$
wybierzmy niezależne zmienne $Z\sim\Normal(0,1)$
oraz $U\sim\chi^2_\nu$ na wspólnej przestrzeni.
Oznaczmy ich rozkłady odpowiednio przez $\mu_Z,\mu_U$.
Nie jest potrzebna jedna przestrzeń dla wszystkich $\nu$.
Z własności rozkładu chi-kwadrat mamy $\E(U/\nu)=1$ i
$\Var(U/\nu)=2/\nu$. Nierówność Czebyszewa daje
\[
\Prob(|U/\nu-1|\geq\varepsilon)\leq
\frac{2}{\nu\varepsilon^2}.
\]
Dla $u>0$ nierówność $z/\sqrt{u/\nu}\leq x$
jest równoważna $z\leq x\sqrt{u/\nu}$.
Z niezależności $Z$ i $U$ oraz twierdzenia Tonellego
dla nieujemnej funkcji charakterystycznej zbioru
$\{(z,u):u>0,\ z\leq x\sqrt{u/\nu}\}$ otrzymujemy
\[
F_{t_\nu}(x)=
\int_{(0,\infty)}\int_{(-\infty,x\sqrt{u/\nu}]}
\dd\mu_Z(z)\,\dd\mu_U(u)
=\E\Phi(x\sqrt{U/\nu}),
\qquad x\in\R,
\]
gdzie $\Phi$ jest dystrybuantą standardowego rozkładu
normalnego, a wartość podcałkową na $\{U\leq0\}$
można określić dowolnie w $[0,1]$.
Dla ustalonego $x\in\R$ i $\eta>0$ z ciągłości funkcji
$v\mapsto\Phi(x\sqrt v)$ w $1$ wybieramy
$\varepsilon\in(0,1)$ tak, by dla $|v-1|<\varepsilon$
różnica $|\Phi(x\sqrt v)-\Phi(x)|$ była mniejsza od $\eta$.
Rozdzielając całkę na zdarzenia
$\{|U/\nu-1|<\varepsilon\}$ i jego dopełnienie, mamy
\[
|F_{t_\nu}(x)-\Phi(x)|
\leq\eta+\frac{2}{\nu\varepsilon^2}.
\]
Najpierw przechodzimy do granicy górnej przy
$\nu\to\infty$, a następnie do granicy $\eta\downarrow0$.
Otrzymujemy zbieżność dystrybuant w każdym punkcie,
czyli zbieżność według rozkładu
(rozdział~\ref{ch:convergence}). Dowód nie wymaga
całkowitości liczby stopni swobody.
\end{uwaga}

\section{Rozkład Cauchy'ego}
Symetryczna, gładka gęstość nie musi opisywać rozkładu z dobrze
określoną średnią. Rozkład Cauchy'ego stanowi ważny kontrprzykład:
mimo symetrii jego dodatnia i ujemna część pierwszego momentu
są rozbieżne. Na nim zobaczymy, dlaczego formalne „znoszenie się”
nieskończoności nie jest dopuszczalnym sposobem obliczania
wartości oczekiwanej.

\begin{definicja}
Niech $x_0\in\R$ i $\gamma>0$. Rozkładem Cauchy'ego
o parametrze położenia $x_0$ i skali $\gamma$
nazywamy rozkład o gęstości
\begin{equation}\label{eq:cauchy}
f(x)=\frac1{\pi\gamma}
\frac{1}{1+((x-x_0)/\gamma)^2},\qquad x\in\R.
\end{equation}
\end{definicja}

\begin{stwierdzenie}
Jeżeli $X$ ma powyższy rozkład, to
\[
F_X(x)=\frac12+\frac1\pi
\arctan\!\left(\frac{x-x_0}{\gamma}\right).
\]
Dla każdego rzeczywistego $r>0$ zachodzi
$\E|X-x_0|^r<\infty$ wtedy i tylko wtedy, gdy $r<1$.
W szczególności wartość oczekiwana $\E X$ nie istnieje,
a wariancja nie jest zdefiniowana.
Rozkład standardowy Cauchy'ego ($x_0=0$, $\gamma=1$)
jest równy rozkładowi $t_1$.
\end{stwierdzenie}
\begin{proof}
Gęstość \eqref{eq:cauchy} jest ciągła i nieujemna.
Podstawiając $v=(u-x_0)/\gamma$, otrzymujemy
\[
\int_{-\infty}^x f(u)\,\dd\lambda_1(u)
=\frac1\pi\int_{-\infty}^{(x-x_0)/\gamma}
\frac{1}{1+v^2}\,\dd\lambda_1(v)
=\frac12+\frac1\pi\arctan\!\left(\frac{x-x_0}{\gamma}\right).
\]
Przy $x\to\infty$ prawa strona dąży do $1$, co
uzasadnia normalizację gęstości i wzór na dystrybuantę.
Dla $r>0$, stosując \eqref{eq:lotus} i to samo
podstawienie, mamy
\[
\E|X-x_0|^r=\frac{2\gamma^r}{\pi}
\int_0^\infty\frac{v^r}{1+v^2}\,\dd\lambda_1(v).
\]
Całka na $(0,1)$ jest skończona, gdyż funkcja
podcałkowa jest ograniczona. Dla $v\geq1$ zachodzi
$\tfrac12 v^{r-2}\leq v^r/(1+v^2)\leq v^{r-2}$.
Całka jest więc skończona dokładnie wtedy, gdy $r<1$.
W szczególności z symetrii rozkładu $X-x_0$ obie części
$(X-x_0)^+$ i $(X-x_0)^-$ mają nieskończone wartości
oczekiwane. Funkcje $x\mapsto x^+$ i $x\mapsto x^-$
są lipschitzowskie ze stałą $1$, stąd
$(X-x_0)^\pm\leq X^\pm+|x_0|$.
Z monotoniczności całki nieujemnej obie wartości
$\E X^+$ i $\E X^-$ również są nieskończone.
Nie istnieje więc $\E X$, także w sensie rozszerzonym.
Podstawienie $\nu=1$ do \eqref{eq:student} i tożsamości
$\Gamma(1)=1$, $\Gamma(1/2)=\sqrt\pi$ daje gęstość
$1/[\pi(1+x^2)]$ rozkładu standardowego Cauchy'ego.
\end{proof}

\begin{przyklad}[Wartość główna a wartość oczekiwana]
Dla standardowego rozkładu Cauchy'ego zachodzi
$\int_{-R}^{R}x/[\pi(1+x^2)]\,\dd\lambda_1(x)=0$
dla każdego $R>0$ z powodu nieparzystości funkcji.
Wartość główna Cauchy'ego całki wynosi zatem zero,
ale nie jest to wartość oczekiwana: całki z części dodatniej
i ujemnej są obie nieskończone. Nie można stosować
symetrycznego znoszenia się rozbieżnych całek w definicji
wartości oczekiwanej.
\end{przyklad}

\begin{przyklad}[Tangens kąta jednostajnego]
Niech $U\sim\Unif(-\pi/2,\pi/2)$. Definiujemy
$X=\tan U$ na zdarzeniu $\{-\pi/2<U<\pi/2\}$,
a poza nim $X=0$. Dopełnienie tego zdarzenia jest
mierzalnym zbiorem zerowym, a tak określona zmienna
jest mierzalna. Funkcja $\tan$ jest dyfeomorfizmem
klasy $C^1$ z przedziału $(-\pi/2,\pi/2)$ na $\R$,
a pochodna funkcji odwrotnej wynosi
$(\arctan x)'=1/(1+x^2)>0$.
Wzór na transformację gęstości z rozdziału~\ref{ch:random-variables}
daje $f_X(x)=1/[\pi(1+x^2)]$, czyli gęstość standardowego
rozkładu Cauchy'ego.
\end{przyklad}


\begin{przyklad}[Porównanie ogonów rozkładu normalnego i Cauchy'ego]
Niech $Z\sim\Normal(0,1)$, $C$ ma standardowy rozkład
Cauchy'ego i niech $t>0$. Ponieważ $\{|Z|\geq t\}=\{Z^2\geq t^2\}$,
z nierówności Markowa dla $Z^2$ oraz $\E Z^2=1$
wynika $\Prob(|Z|\geq t)\leq1/t^2$.
W rzeczywistości ogon normalny maleje znacznie szybciej:
dla każdego $u>0$, korzystając z
$M_Z(u)=e^{u^2/2}$ i stosując nierówność Markowa
do nieujemnej zmiennej $e^{uZ}$ z progiem $e^{ut}>0$, mamy
$\Prob(Z\geq t)=\Prob(e^{uZ}\geq e^{ut})
\leq e^{-ut+u^2/2}$.
Wybór $u=t$ daje $\Prob(Z\geq t)\leq e^{-t^2/2}$,
a przez symetrię $\Prob(|Z|\geq t)\leq2e^{-t^2/2}$.

Tymczasem dla $C$ zachodzi
$\Prob(|C|\geq t)=1-(2/\pi)\arctan t$,
gdyż rozkład jest symetryczny i nie ma mas punktowych.
Dla $t>0$ mamy $\pi/2-\arctan t=\arctan(1/t)$,
a z różniczkowalności funkcji $\arctan$ w zerze
wynika $\arctan v/v\to1$ przy $v\downarrow0$.
W konsekwencji $\Prob(|C|\geq t)\sim 2/(\pi t)$
przy $t\to\infty$.
Dla nieujemnej zmiennej losowej $Y$ na przestrzeni
$(\Omega,\mathcal F,\Prob)$ rozważmy zbiór
$D=\{(\omega,t)\in\Omega\times(0,\infty):t<Y(\omega)\}$.
Jest on mierzalny w produktowej $\sigma$-algebrze,
ponieważ
\[
D=\bigcup_{\substack{q\in\mathbb Q\\q>0}}
\bigl(\{Y>q\}\times(0,q)\bigr).
\]
Dla każdego $\omega$ przekrój zbioru $D$ wzdłuż
drugiej współrzędnej jest przedziałem $(0,Y(\omega))$.
Zatem, stosując \atlasref{ME-12::thm:tonelli-12}{twierdzenie Tonellego}
do nieujemnej funkcji $\chi_D$ względem miar
$\Prob$ i $\lambda_1|_{(0,\infty)}$,
otrzymujemy (również dla nieskończonej całki)
\[
\begin{aligned}
\E Y
&=\int_\Omega\int_0^\infty\chi_D(\omega,t)\,
\dd\lambda_1(t)\,\dd\Prob(\omega)\\
&=\int_0^\infty\int_\Omega\chi_D(\omega,t)\,
\dd\Prob(\omega)\,\dd\lambda_1(t)
=\int_0^\infty\Prob(Y>t)\,\dd\lambda_1(t).
\end{aligned}
\]
Oszacowanie ogona $|C|$ od dołu przez dodatnią stałą
razy $1/t$ dla dużych $t$ daje $\E|C|=\infty$,
zgodnie z wcześniejszym rachunkiem momentów.
Rysunek~\ref{fig:rp1-06-continuous-densities} przedstawia porównanie trzech symetrycznych gęstości.
\end{przyklad}

\begin{figure}[htbp]
\centering
\begin{tikzpicture}[x=.79cm,y=10.4cm,>=Stealth]
\draw[->] (-7.15,0)--(7.5,0) node[right] {$x$};
\draw[->] (0,0)--(0,.445) node[above] {$f(x)$};
\foreach \x in {-6,-4,-2,2,4,6} \node[below,font=\small] at (\x,0) {$\x$};
\foreach \y/\lab in {.1/{0{,}1},.2/{0{,}2},.3/{0{,}3},.4/{0{,}4}} {
 \draw[gray!20] (-7,\y)--(7,\y);
 \node[left,font=\scriptsize] at (0,\y) {$\lab$};
}
\draw[blue!75!black,very thick] plot coordinates {(-7,0) (-6.9,0) (-6.8,0) (-6.7,0) (-6.6,0) (-6.5,0) (-6.4,0) (-6.3,0) (-6.2,0) (-6.1,0) (-6,0) (-5.9,0) (-5.8,0) (-5.7,0) (-5.6,1e-7) (-5.5,1e-7) (-5.4,2e-7) (-5.3,3e-7) (-5.2,5e-7) (-5.1,9e-7) (-5,0.0000015) (-4.9,0.0000024) (-4.8,0.000004) (-4.7,0.0000064) (-4.6,0.0000101) (-4.5,0.000016) (-4.4,0.0000249) (-4.3,0.0000385) (-4.2,0.0000589) (-4.1,0.0000893) (-4,0.0001338) (-3.9,0.0001987) (-3.8,0.0002919) (-3.7,0.0004248) (-3.6,0.0006119) (-3.5,0.0008727) (-3.4,0.0012322) (-3.3,0.0017226) (-3.2,0.0023841) (-3.1,0.0032668) (-3,0.0044318) (-2.9,0.0059525) (-2.8,0.0079155) (-2.7,0.0104209) (-2.6,0.013583) (-2.5,0.0175283) (-2.4,0.0223945) (-2.3,0.028327) (-2.2,0.0354746) (-2.1,0.0439836) (-2,0.053991) (-1.9,0.0656158) (-1.8,0.0789502) (-1.7,0.0940491) (-1.6,0.1109208) (-1.5,0.1295176) (-1.4,0.1497275) (-1.3,0.1713686) (-1.2,0.1941861) (-1.1,0.2178522) (-1,0.2419707) (-0.9,0.2660852) (-0.8,0.2896916) (-0.7,0.3122539) (-0.6,0.3332246) (-0.5,0.3520653) (-0.4,0.3682701) (-0.3,0.3813878) (-0.2,0.3910427) (-0.1,0.3969525) (0,0.3989423) (0.1,0.3969525) (0.2,0.3910427) (0.3,0.3813878) (0.4,0.3682701) (0.5,0.3520653) (0.6,0.3332246) (0.7,0.3122539) (0.8,0.2896916) (0.9,0.2660852) (1,0.2419707) (1.1,0.2178522) (1.2,0.1941861) (1.3,0.1713686) (1.4,0.1497275) (1.5,0.1295176) (1.6,0.1109208) (1.7,0.0940491) (1.8,0.0789502) (1.9,0.0656158) (2,0.053991) (2.1,0.0439836) (2.2,0.0354746) (2.3,0.028327) (2.4,0.0223945) (2.5,0.0175283) (2.6,0.013583) (2.7,0.0104209) (2.8,0.0079155) (2.9,0.0059525) (3,0.0044318) (3.1,0.0032668) (3.2,0.0023841) (3.3,0.0017226) (3.4,0.0012322) (3.5,0.0008727) (3.6,0.0006119) (3.7,0.0004248) (3.8,0.0002919) (3.9,0.0001987) (4,0.0001338) (4.1,0.0000893) (4.2,0.0000589) (4.3,0.0000385) (4.4,0.0000249) (4.5,0.000016) (4.6,0.0000101) (4.7,0.0000064) (4.8,0.000004) (4.9,0.0000024) (5,0.0000015) (5.1,9e-7) (5.2,5e-7) (5.3,3e-7) (5.4,2e-7) (5.5,1e-7) (5.6,1e-7) (5.7,0) (5.8,0) (5.9,0) (6,0) (6.1,0) (6.2,0) (6.3,0) (6.4,0) (6.5,0) (6.6,0) (6.7,0) (6.8,0) (6.9,0) (7,0)};
\draw[teal!70!black,dashed,very thick] plot coordinates {(-7,0.0012234) (-6.9,0.0012915) (-6.8,0.0013644) (-6.7,0.0014424) (-6.6,0.0015259) (-6.5,0.0016156) (-6.4,0.0017118) (-6.3,0.0018151) (-6.2,0.0019263) (-6.1,0.0020459) (-6,0.0021749) (-5.9,0.0023139) (-5.8,0.0024641) (-5.7,0.0026263) (-5.6,0.0028019) (-5.5,0.0029921) (-5.4,0.0031984) (-5.3,0.0034223) (-5.2,0.0036657) (-5.1,0.0039307) (-5,0.0042194) (-4.9,0.0045343) (-4.8,0.0048784) (-4.7,0.0052549) (-4.6,0.0056672) (-4.5,0.0061195) (-4.4,0.0066163) (-4.3,0.0071629) (-4.2,0.007765) (-4.1,0.0084293) (-4,0.0091634) (-3.9,0.0099757) (-3.8,0.010876) (-3.7,0.0118754) (-3.6,0.0129866) (-3.5,0.014224) (-3.4,0.0156041) (-3.3,0.0171458) (-3.2,0.0188706) (-3.1,0.0208033) (-3,0.022972) (-2.9,0.0254092) (-2.8,0.0281516) (-2.7,0.0312415) (-2.6,0.0347266) (-2.5,0.0386615) (-2.4,0.0431076) (-2.3,0.0481341) (-2.2,0.0538183) (-2.1,0.0602456) (-2,0.0675097) (-1.9,0.075711) (-1.8,0.0849558) (-1.7,0.0953524) (-1.6,0.1070071) (-1.5,0.1200172) (-1.4,0.1344617) (-1.3,0.1503891) (-1.2,0.1678016) (-1.1,0.186637) (-1,0.2067483) (-0.9,0.2278831) (-0.8,0.2496659) (-0.7,0.2715884) (-0.6,0.2930107) (-0.5,0.3131809) (-0.4,0.3312744) (-0.3,0.3464536) (-0.2,0.3579438) (-0.1,0.3651144) (0,0.3675526) (0.1,0.3651144) (0.2,0.3579438) (0.3,0.3464536) (0.4,0.3312744) (0.5,0.3131809) (0.6,0.2930107) (0.7,0.2715884) (0.8,0.2496659) (0.9,0.2278831) (1,0.2067483) (1.1,0.186637) (1.2,0.1678016) (1.3,0.1503891) (1.4,0.1344617) (1.5,0.1200172) (1.6,0.1070071) (1.7,0.0953524) (1.8,0.0849558) (1.9,0.075711) (2,0.0675097) (2.1,0.0602456) (2.2,0.0538183) (2.3,0.0481341) (2.4,0.0431076) (2.5,0.0386615) (2.6,0.0347266) (2.7,0.0312415) (2.8,0.0281516) (2.9,0.0254092) (3,0.022972) (3.1,0.0208033) (3.2,0.0188706) (3.3,0.0171458) (3.4,0.0156041) (3.5,0.014224) (3.6,0.0129866) (3.7,0.0118754) (3.8,0.010876) (3.9,0.0099757) (4,0.0091634) (4.1,0.0084293) (4.2,0.007765) (4.3,0.0071629) (4.4,0.0066163) (4.5,0.0061195) (4.6,0.0056672) (4.7,0.0052549) (4.8,0.0048784) (4.9,0.0045343) (5,0.0042194) (5.1,0.0039307) (5.2,0.0036657) (5.3,0.0034223) (5.4,0.0031984) (5.5,0.0029921) (5.6,0.0028019) (5.7,0.0026263) (5.8,0.0024641) (5.9,0.0023139) (6,0.0021749) (6.1,0.0020459) (6.2,0.0019263) (6.3,0.0018151) (6.4,0.0017118) (6.5,0.0016156) (6.6,0.0015259) (6.7,0.0014424) (6.8,0.0013644) (6.9,0.0012915) (7,0.0012234)};
\draw[orange!85!black,dashdotted,very thick] plot coordinates {(-7,0.0063662) (-6.9,0.0065482) (-6.8,0.0067381) (-6.7,0.0069364) (-6.6,0.0071434) (-6.5,0.0073598) (-6.4,0.007586) (-6.3,0.0078228) (-6.2,0.0080707) (-6.1,0.0083305) (-6,0.008603) (-5.9,0.0088889) (-5.8,0.0091891) (-5.7,0.0095046) (-5.6,0.0098365) (-5.5,0.0101859) (-5.4,0.010554) (-5.3,0.0109422) (-5.2,0.011352) (-5.1,0.0117849) (-5,0.0122427) (-4.9,0.0127273) (-4.8,0.0132408) (-4.7,0.0137856) (-4.6,0.0143642) (-4.5,0.0149793) (-4.4,0.0156341) (-4.3,0.016332) (-4.2,0.0170767) (-4.1,0.0178725) (-4,0.0187241) (-3.9,0.0196366) (-3.8,0.0206159) (-3.7,0.0216685) (-3.6,0.0228016) (-3.5,0.0240234) (-3.4,0.0253431) (-3.3,0.0267712) (-3.2,0.0283194) (-3.1,0.0300009) (-3,0.031831) (-2.9,0.0338268) (-2.8,0.0360079) (-2.7,0.0383968) (-2.6,0.0410193) (-2.5,0.0439048) (-2.4,0.0470873) (-2.3,0.0506057) (-2.2,0.0545051) (-2.1,0.0588373) (-2,0.063662) (-1.9,0.0690477) (-1.8,0.0750731) (-1.7,0.0818277) (-1.6,0.0894129) (-1.5,0.0979415) (-1.4,0.1075371) (-1.3,0.1183308) (-1.2,0.1304549) (-1.1,0.1440316) (-1,0.1591549) (-0.9,0.1758618) (-0.8,0.1940914) (-0.7,0.2136308) (-0.6,0.2340514) (-0.5,0.2546479) (-0.4,0.2744051) (-0.3,0.2920274) (-0.2,0.3060672) (-0.1,0.3151583) (0,0.3183099) (0.1,0.3151583) (0.2,0.3060672) (0.3,0.2920274) (0.4,0.2744051) (0.5,0.2546479) (0.6,0.2340514) (0.7,0.2136308) (0.8,0.1940914) (0.9,0.1758618) (1,0.1591549) (1.1,0.1440316) (1.2,0.1304549) (1.3,0.1183308) (1.4,0.1075371) (1.5,0.0979415) (1.6,0.0894129) (1.7,0.0818277) (1.8,0.0750731) (1.9,0.0690477) (2,0.063662) (2.1,0.0588373) (2.2,0.0545051) (2.3,0.0506057) (2.4,0.0470873) (2.5,0.0439048) (2.6,0.0410193) (2.7,0.0383968) (2.8,0.0360079) (2.9,0.0338268) (3,0.031831) (3.1,0.0300009) (3.2,0.0283194) (3.3,0.0267712) (3.4,0.0253431) (3.5,0.0240234) (3.6,0.0228016) (3.7,0.0216685) (3.8,0.0206159) (3.9,0.0196366) (4,0.0187241) (4.1,0.0178725) (4.2,0.0170767) (4.3,0.016332) (4.4,0.0156341) (4.5,0.0149793) (4.6,0.0143642) (4.7,0.0137856) (4.8,0.0132408) (4.9,0.0127273) (5,0.0122427) (5.1,0.0117849) (5.2,0.011352) (5.3,0.0109422) (5.4,0.010554) (5.5,0.0101859) (5.6,0.0098365) (5.7,0.0095046) (5.8,0.0091891) (5.9,0.0088889) (6,0.008603) (6.1,0.0083305) (6.2,0.0080707) (6.3,0.0078228) (6.4,0.007586) (6.5,0.0073598) (6.6,0.0071434) (6.7,0.0069364) (6.8,0.0067381) (6.9,0.0065482) (7,0.0063662)};
\draw[blue!75!black,very thick] (3.25,.39)--(4,.39) node[right,font=\small] {$\mathcal N(0,1)$};
\draw[teal!70!black,very thick,dashed] (3.25,.34)--(4,.34) node[right,font=\small] {$t_3$};
\draw[orange!85!black,very thick,dashdotted] (3.25,.29)--(4,.29) node[right,font=\small] {Cauchy};
\end{tikzpicture}
\caption{Gęstości standardowego rozkładu normalnego,
rozkładu Studenta z trzema stopniami swobody oraz
standardowego rozkładu Cauchy'ego ($t_1$).
Wzory na gęstości pochodzą z wcześniejszych sekcji.
Cięższy ogon nie oznacza większej gęstości w każdym punkcie.}
\label{fig:rp1-06-continuous-densities}
\end{figure}


\section{Dalsze użyteczne rozkłady}
Nie wszystkie użyteczne rodziny wymagają nowej konstrukcji od zera.
Rozkład logarytmiczno-normalny powstaje przez wykładnicze przekształcenie
rozkładu normalnego. Rozkład Fishera--Snedecora powstaje z ilorazu
dwóch niezależnych odpowiednio unormowanych zmiennych chi-kwadrat.
To przykłady ważnej zasady: z prostych rozkładów tworzymy
bardziej złożone za pomocą działań na zmiennych losowych.

\begin{definicja}
Zmienna losowa $X$ ma rozkład logarytmiczno-normalny
z parametrami $\mu\in\R$, $\sigma>0$, jeżeli $\Prob(X>0)=1$
oraz $\log X\sim\Normal(\mu,\sigma^2)$, przy czym
$\log X$ określamy dowolnie w sposób mierzalny na
zdarzeniu $\{X\leq0\}$ o prawdopodobieństwie zero.
Równoważnie jest to rozkład zmiennej $e^Y$, gdzie
$Y\sim\Normal(\mu,\sigma^2)$: w pierwszym kierunku
$X=e^{\log X}$ prawie na pewno, a w drugim
$\log(e^Y)=Y$ punktowo.
Z twierdzenia o transformacji monotonicznej
(rozdział~\ref{ch:random-variables}), zastosowanego do
dyfeomorfizmu $\exp:\R\to(0,\infty)$ o odwrotności
$\log$ i pochodnej odwrotności $1/x$, wynika gęstość
względem miary Lebesgue'a:
\[
f_X(x)=
\begin{cases}
\displaystyle\frac{\exp(-(\log x-\mu)^2/(2\sigma^2))}
{x\sigma\sqrt{2\pi}},&x>0,\\
0,&x\leq0.
\end{cases}
\]
Jej całka wynosi jeden na mocy podstawienia $y=\log x$
i normalizacji gęstości normalnej.
\end{definicja}

\begin{stwierdzenie}
Jeżeli $X$ ma rozkład logarytmiczno-normalny
z parametrami $\mu,\sigma$, to dla każdego $r\in\R$
zachodzi $\E X^r=\exp(r\mu+r^2\sigma^2/2)$,
gdzie potęgi określamy mierzalnie i nieujemnie
na zbiorze zerowym $\{X\leq0\}$.
W szczególności $\E X=e^{\mu+\sigma^2/2}$ oraz
$\Var(X)=(e^{\sigma^2}-1)e^{2\mu+\sigma^2}$.
\end{stwierdzenie}
\begin{proof}
Mamy $X=e^Y$ prawie na pewno dla
$Y\sim\Normal(\mu,\sigma^2)$.
Stąd $\E X^r=\E e^{rY}=M_Y(r)=\exp(r\mu+r^2\sigma^2/2)$
dla dowolnego rzeczywistego $r$, ponieważ funkcja tworząca
momenty rozkładu normalnego jest skończona na całej prostej.
Dla $r=1,2$ daje to odpowiednio $e^{\mu+\sigma^2/2}$
i $e^{2\mu+2\sigma^2}$, a więc
\[
\Var(X)=e^{2\mu+2\sigma^2}-e^{2\mu+\sigma^2}
=(e^{\sigma^2}-1)e^{2\mu+\sigma^2}.
\]
\end{proof}
\begin{uwaga}
Nie należy stąd wnioskować, że funkcja tworząca momenty
$M_X(t)=\E e^{tX}$ jest skończona w otoczeniu zera.
Dla $t\leq0$ mamy $0<e^{tX}\leq1$ prawie na pewno,
więc $M_X(t)<\infty$. Natomiast dla $t>0$, stosując
\eqref{eq:lotus} do $Y$ i nieujemnej funkcji
$y\mapsto e^{te^y}$, otrzymujemy
\[
M_X(t)=\frac{1}{\sigma\sqrt{2\pi}}
\int_\R\exp\!\left(te^y-\frac{(y-\mu)^2}{2\sigma^2}\right)
\,\dd\lambda_1(y).
\]
Ponieważ $(y-\mu)^2e^{-y}\to0$ przy $y\to+\infty$,
wykładnik można zapisać jako
$e^y(t-(y-\mu)^2e^{-y}/(2\sigma^2))$;
dla dostatecznie dużego $y$ jest on nie mniejszy
niż $(t/2)e^y$ i dąży do $+\infty$.
Zatem funkcja podcałkowa jest od pewnego miejsca
nie mniejsza niż $1$. Całka jest nieskończona.
Dziedziną skończoności $M_X$ jest zatem $(-\infty,0]$,
mimo istnienia wszystkich dodatnich momentów.
Zob. również kontrprzykład dotyczący jednoznaczności
rozkładu wyznaczonego przez momenty w rozdziale~\ref{ch:expectation}.
\end{uwaga}

\begin{definicja}
Niech $d_1,d_2\geq1$ będą liczbami całkowitymi, a
$U\sim\chi^2_{d_1}$ oraz $V\sim\chi^2_{d_2}$ niezależnymi
zmiennymi losowymi na wspólnej przestrzeni probabilistycznej.
Ponieważ $\Prob(U>0,V>0)=1$, definiujemy
\[
F=\begin{cases}
(U/d_1)/(V/d_2),&U>0,\ V>0,\\
0,&\text{w przeciwnym przypadku}.
\end{cases}
\]
Zmienna jest mierzalna, gdyż określająca ją funkcja
jest ciągła na $(0,\infty)^2$ i przedłużona zerem
poza tym obszarem. Jej rozkład zależy wyłącznie od
$d_1,d_2$, ponieważ jest obrazem ustalonego produktowego
rozkładu pary $(U,V)$. Nazywamy go rozkładem
Fishera--Snedecora i oznaczamy $F_{d_1,d_2}$.
\end{definicja}

\begin{stwierdzenie}
Rozkład $F_{d_1,d_2}$ ma gęstość
\[
f(x)=
\begin{cases}
\displaystyle\frac{\Gamma((d_1+d_2)/2)}
{\Gamma(d_1/2)\Gamma(d_2/2)}
\left(\frac{d_1}{d_2}\right)^{d_1/2}
\frac{x^{d_1/2-1}}{(1+d_1x/d_2)^{(d_1+d_2)/2}},
&x>0,\\
0,&x\leq0.
\end{cases}
\]
Jeżeli $T\sim t_\nu$ i $\nu\geq1$ jest liczbą całkowitą, to $T^2\sim F_{1,\nu}$.
\end{stwierdzenie}
\begin{proof}
Rozpatrzmy niezależne $U,V$ o rozkładach chi-kwadrat
z definicji $F_{d_1,d_2}$.
Oznaczmy $a=d_1/2$ oraz $b=d_2/2$.
Z niezależności $U,V$ oraz wzoru na gęstość rozkładu
łącznego (rozdział~\ref{ch:random-variables}) ich gęstość
łączna na $(0,\infty)^2$ jest równa
\[
f_{U,V}(u,v)=\frac{u^{a-1}v^{b-1}e^{-(u+v)/2}}
{2^{a+b}\Gamma(a)\Gamma(b)}.
\]
Oznaczmy $X=F$. Na dodatnim kwadrancie odwzorowanie
$(u,v)\mapsto((d_2/d_1)u/v,v)$ jest dyfeomorfizmem
klasy $C^1$ na siebie. Jego odwrotność ma postać
$(x,v)\mapsto((d_1/d_2)xv,v)$, a macierz jej pochodnych
wynosi
\[
\begin{pmatrix}
(d_1/d_2)v&(d_1/d_2)x\\
0&1
\end{pmatrix}.
\]
Moduł wyznacznika jest równy $(d_1/d_2)v>0$.
Para $(U,V)$ jest skupiona na dodatnim kwadrancie,
więc \atlasref{ME-13::thm:change-variables-13}{twierdzenie o zamianie zmiennych}
daje gęstość pary $(X,V)$ przez podstawienie do gęstości
łącznej i pomnożenie przez moduł jakobianu.
Stosując \atlasref{ME-12::thm:tonelli-12}{twierdzenie Tonellego}
do nieujemnej gęstości, a następnie całkując po $v>0$,
dostajemy
\[
f_X(x)=\frac{(d_1/d_2)^a x^{a-1}}
{2^{a+b}\Gamma(a)\Gamma(b)}
\int_0^\infty v^{a+b-1}
\exp\!\left(-\frac{1+(d_1/d_2)x}{2}v\right)
\,\dd\lambda_1(v),\quad x>0.
\]
Dla $x>0$ liczba $c=(1+(d_1/d_2)x)/2$ jest dodatnia.
Po podstawieniu $w=cv$ całka wynosi
$c^{-(a+b)}\Gamma(a+b)$, ponieważ $a+b>0$.
Po uproszczeniu otrzymujemy podaną gęstość dla $x>0$;
dla $x\leq0$ przyjmujemy gęstość równą zeru,
gdyż $X>0$ prawie na pewno.
Całka z uzyskanej gęstości jest równa jeden: jest to
gęstość brzegowa rozkładu probabilistycznego $(X,V)$,
a twierdzenie Tonellego pozwala zintegrować ją po $\R$.

Jeżeli $T\sim t_\nu$, to z definicji rozkładu
$\Law(T)=\Law(Z/\sqrt{V/\nu})$, gdzie
$Z\sim\Normal(0,1)$, $V\sim\chi_\nu^2$ są niezależne,
a iloraz określamy zerem na $\{V\leq0\}$.
Ponieważ rozkład mierzalnego przekształcenia zależy
tylko od rozkładu jego argumentu, mamy
$\Law(T^2)=\Law((Z^2/1)/(V/\nu))$, z dokładnością
do modyfikacji ilorazu na zbiorze zerowym.
Z poprzedniego twierdzenia o sumie kwadratów wynika
$Z^2\sim\chi_1^2$, a $Z^2,V$ są niezależne jako
mierzalne funkcje niezależnych zmiennych.
Zatem $T^2\sim F_{1,\nu}$. Ten argument działa
również dla dowolnego rzeczywistego $\nu>0$, jeśli
rozszerzymy definicję $F_{d_1,d_2}$ na rzeczywiste
dodatnie parametry; w tekście zasadnicza definicja
podana została dla całkowitych stopni swobody.
\end{proof}

\section{Porównanie rozkładów i wybór modelu}
Właściwy wybór modelu nie polega na dopasowaniu znanego wzoru
do przypadkowej tabeli. Najpierw trzeba rozpoznać mechanizm:
losowanie ze zwracaniem czy bez, ustalona liczba prób czy czekanie
na sukces, pojedyncza wielkość czy suma/iloraz wielkości losowych.
Dopiero potem parametryzuje się odpowiadający mu rozkład.

\begin{center}
\small
\begin{tabular}{@{}p{3.0cm}p{4.1cm}p{5.9cm}@{}}
\hline
Rozkład & Parametry & Typowe zastosowanie \\
\hline
Bernoulliego & $p\in[0,1]$ & Wynik pojedynczej próby binarnej \\
Dwumianowy & $n\in\N_0$, $p\in[0,1]$ & Liczba sukcesów w niezależnych próbach \\
Hipergeometryczny & $N\geq1$, $0\leq K,n\leq N$ (całkowite) & Liczba sukcesów przy losowaniu bez zwracania \\
Poissona & $\lambda>0$ & Przybliżenie liczby sukcesów w wielu niezależnych próbach o małym prawdopodobieństwie sukcesu \\
Geometryczny & $p\in(0,1]$ & Liczba prób do pierwszego sukcesu \\
Ujemny dwumianowy & $r\geq1$ całkowite, $p\in(0,1]$ & Liczba prób do $r$-tego sukcesu \\
Jednostajny & $a<b$ & Jednorodne losowanie punktu odcinka \\
Wykładniczy & $\lambda>0$ & Niepamiętający czas oczekiwania \\
Normalny & $\mu\in\R$, $\sigma>0$ & Błędy pomiaru, sumy licznych składników \\
Gamma & $\alpha,\lambda>0$ & Dodatnie czasy i sumy czasów oczekiwania \\
Chi-kwadrat & $\nu>0$ & Suma kwadratów niezależnych standardowych zmiennych normalnych dla całkowitego $\nu$ \\
Beta & $a,b>0$ & Udziały, proporcje i zmienne o wartościach w $(0,1)$ \\
t-Studenta & $\nu>0$ & Iloraz standardowej zmiennej normalnej i pierwiastka z niezależnej unormowanej zmiennej chi-kwadrat \\
Cauchy'ego & $x_0\in\R$, $\gamma>0$ & Model z ciężkimi ogonami, kontrprzykłady dla momentów \\
Logarytmiczno-normalny & $\mu\in\R$, $\sigma>0$ & Dodatnia zmienna o normalnym logarytmie \\
Fishera--Snedecora & $d_1,d_2\geq1$ całkowite & Iloraz niezależnych unormowanych zmiennych chi-kwadrat \\
\hline
\end{tabular}
\end{center}

\begin{uwaga}
Zapis $\N_0$ w tabeli oznacza $\{0,1,2,\ldots\}$.
Przykłady zastosowań są wskazówkami modelowymi, a nie
warunkami definiującymi rozkłady. To, czy model jest odpowiedni,
wynika z przyjętych założeń, sposobu zbierania danych
i weryfikacji ich zgodności z modelem.
\end{uwaga}

